\documentclass[11pt,a4paper]{article}

\input{glyphtounicode}
\pdfgentounicode=1

\usepackage{cmap}
\usepackage[utf8]{inputenc}
\usepackage[T1]{fontenc}
\usepackage{lmodern}
\usepackage{amsmath,amssymb,amsthm,mathtools}
\usepackage{mathrsfs}
\usepackage{geometry}
\geometry{a4paper,margin=1in}
\usepackage{enumitem}
\usepackage{graphicx}
\usepackage{booktabs}
\usepackage{longtable}
\usepackage{array}
\usepackage{tabularx}
\usepackage{tikz-cd}
\usepackage[backend=biber]{biblatex}
\addbibresource{core/canonical_core_references.bib}
\usepackage[final]{microtype}
\usepackage{xcolor}
\usepackage[unicode,pdfencoding=auto,psdextra,hidelinks]{hyperref}
\hypersetup{
  pdftitle={BaseCore of the Rigid Identity Framework},
  pdfauthor={Johnny Andrey P\'erez Ram\'irez},
  pdfsubject={Autonomous Mathematical Source Text for Foundational Dynamics, Identity Rigidity, Regime Constraints, and Operational-Criterion Grammar},
  pdfkeywords={basecore, rigid identity, inverse limit, contractive dynamics, non-finite determination, operational criterion, claim boundary}
}

\raggedbottom
\widowpenalty=10000
\clubpenalty=10000
\setlength{\emergencystretch}{3em}
\DisableLigatures{encoding = *, family = *}

\theoremstyle{plain}
\newtheorem{theorem}{Theorem}[section]
\newtheorem{lemma}[theorem]{Lemma}
\newtheorem{proposition}[theorem]{Proposition}
\newtheorem{corollary}[theorem]{Corollary}

\theoremstyle{definition}
\newtheorem{definition}[theorem]{Definition}
\newtheorem{axiom}[theorem]{Axiom}
\newtheorem{hypothesis}[theorem]{Hypothesis}
\newtheorem{assumption}[theorem]{Assumption}
\newtheorem{criterion}[theorem]{Criterion}
\newtheorem{prediction}[theorem]{Prediction}

\theoremstyle{remark}
\newtheorem{remark}[theorem]{Remark}
\newtheorem{example}[theorem]{Example}
\newtheorem{caveat}[theorem]{Caveat}
\newtheorem{nontheorem}[theorem]{Non-Theorem}

\newcolumntype{Y}{>{\raggedright\arraybackslash}X}
\newcolumntype{P}[1]{>{\raggedright\arraybackslash}p{#1}}

\newcommand{\Hilb}{\mathcal{H}}
\newcommand{\I}{\mathcal{I}}
\newcommand{\U}{\mathcal{U}}
\newcommand{\A}{\mathcal{A}}
\newcommand{\norm}[1]{\left\lVert #1 \right\rVert}
\newcommand{\inner}[2]{\left\langle #1, #2 \right\rangle}

\newcommand{\E}{\mathcal{E}}
\newcommand{\Obs}{\mathbf{Obs}}
\newcommand{\CCR}{\mathbf{CCR}}
\newcommand{\MO}{M_\Omega}
\newcommand{\dw}{d_w}
\newcommand{\dE}{d_{\mathcal{E}}}
\newcommand{\dI}{d_{\mathcal{I}}}
\newcommand{\nullphi}{\varnothing_\phi}
\newcommand{\Attr}{\mathrm{Attr}}
\newcommand{\Eplus}{\mathcal{E}^+}

\newcommand{\Cop}{\mathrm{Crit}_{\mathrm{op}}}
\newcommand{\Qop}{\mathrm{Part}_{\mathrm{op}}}
\newcommand{\Qualop}{\mathrm{Class}_{\mathrm{op}}}
\newcommand{\Iper}{I_{\mathrm{per}}}
\newcommand{\Iri}{I_{\mathrm{ri}}}
\newcommand{\Iint}{I_{\mathrm{int}}}
\newcommand{\Icont}{I_{\mathrm{cont}}}
\newcommand{\Idiff}{I_{\mathrm{diff}}}
\newcommand{\Ileg}{I_{\mathrm{leg}}}
\newcommand{\Aset}{A_S}
\newcommand{\dist}{\operatorname{dist}}
\newcommand{\Lip}{\operatorname{Lip}}
\newcommand{\Id}{\mathcal{I}}
\newcommand{\Read}{\mathscr{R}}
\newcommand{\Hist}{\mathscr{H}}
\newcommand{\Dec}{\mathscr{D}}
\newcommand{\Class}{\mathscr{Q}}
\newcommand{\Part}{\mathscr{P}}
\newcommand{\nullclass}{\mathbf{0}}
\newcommand{\eps}{\varepsilon}
\newcommand{\Nset}{\mathbb{N}}
\newcommand{\Cert}{\mathrm{Cert}}
\newcommand{\Eq}{\mathrm{Eq}}
\newcommand{\Rupt}{\mathrm{Rupt}}

\title{%
\textbf{BaseCore of the Rigid Identity Framework}\\[0.5em]
\large Autonomous Mathematical Source Text for Foundational Dynamics, Identity Rigidity, Regime Constraints, and Operational-Criterion Grammar
}
\author{Johnny Andrey P\'{e}rez Ram\'{i}rez}
\date{}

\begin{document}

\maketitle

\begin{abstract}
This document is the BaseCore of the Rigid Identity Framework: the autonomous mathematical base layer of the source tree. It selectively normalizes formal material from Papers 1--6 only where theorem ownership belongs in the base layer. Papers 7--9 remain downstream and are not required for the validity of any theorem stated here. BaseCore proves contractive projection dynamics, typed computable witnesses, parameterwise non-collapse under explicit anti-constant assumptions, inverse-limit identity under non-finite-determination assumptions, conditional rigidity under admissible perturbation models, conditional non-simulability of CCR targets by finite simulator classes, operational-criterion grammar, and falsation / claim-boundary ledgers. It does not prove phenomenality, metaphysical subjectivity, human equivalence, runtime instantiation, bridge admissibility, or external adjudication. The LaTeX tree is the source of truth; the PDF is a compiled artifact only.
\end{abstract}

\tableofcontents
\newpage

\part{Foundational Dynamics}
\section{Scope, System Boundary, and Non-Inference Note}

\noindent\textbf{What BaseCore does.}\enspace This document proves the mathematical base results that later layers may use: projection existence and uniqueness, strict contraction, unique fixed points, compact attractor families, parameterwise non-collapse under explicit anti-constant assumptions, inverse-limit identity under non-finite-determination assumptions, conditional rigidity, conditional non-simulability, and operational-criterion grammar.

\noindent\textbf{What BaseCore is not.}\enspace BaseCore is not the whole QICN corpus. It is not a bridge program, not runtime validation, not a phenomenality theory, not a human-equivalence claim, and not an external-adjudication layer.

\noindent\textbf{What should not be inferred.}\enspace Nothing in BaseCore alone licenses claims about consciousness, subjectivity, phenomenality, moral status, human-machine equivalence, bridge admissibility, or deployment readiness. Those are downstream burdens, not consequences of the theorems below.

\noindent\textbf{Source-of-truth rule.}\enspace The LaTeX tree is the source of truth. Compiled PDFs, release bundles, and runtime artifacts are derived surfaces only.

\section{Introduction}

\subsection{Purpose}

BaseCore is the canonical base layer of the Rigid Identity Framework. Its role is narrower and stronger than a broad architectural survey: it owns only the theorem-level mathematics that downstream papers may depend on without making BaseCore depend on them.

\subsection{Methodology}

We use standard tools from:
\begin{itemize}[leftmargin=*]
\item Hilbert-space geometry and metric projections;
\item contraction mappings and fixed-point theory;
\item compactness and continuity arguments;
\item projective systems and inverse limits;
\item assumption-led claim-boundary discipline.
\end{itemize}

\subsection{Scope discipline}

This section is pure mathematics. The symbols ``identity,'' ``regime,'' and ``criterion'' are structural labels only. Their ordinary-language interpretations are not part of any proof unless explicitly formalized as assumptions.

\section{Minimal Base Hypotheses}

\begin{hypothesis}[H1: Projection Structure]\label{hyp:H1}
Let $\Hilb$ be a real Hilbert space with inner product $\inner{\cdot}{\cdot}$ and norm $\norm{\cdot}$. Let $\I \subset \Hilb$ be non-empty, closed, and convex.
\end{hypothesis}

\begin{hypothesis}[H2: Contraction Operator]\label{hyp:H2}
Let $K:\Hilb\to\Hilb$ be a bounded linear operator with operator norm $\norm{K}<1$.
\end{hypothesis}

\begin{hypothesis}[H3: Completeness]\label{hyp:H3}
The ambient space $\Hilb$ is complete.
\end{hypothesis}

\begin{hypothesis}[H4: Parameter Compactness]\label{hyp:H4}
Let $\U$ be a compact metric space and let $\Gamma:\U\to\Hilb$ be uniformly continuous.
\end{hypothesis}

\begin{proposition}[Base sufficiency of H1--H4]\label{prop:minimal}
Hypotheses H1--H4 are sufficient for:
\begin{enumerate}[leftmargin=*]
\item existence and uniqueness of the metric projection onto $\I$;
\item strict contraction of the transition operators $T_u$;
\item existence and uniqueness of the fixed point $f_u^\ast$ for each $u\in\U$;
\item compactness of the attractor family $\A^\ast=\{f_u^\ast:u\in\U\}$.
\end{enumerate}
They do not, by themselves, imply parameterwise non-collapse.
\end{proposition}

\begin{proof}
The four items are exactly the content of Theorems~\ref{thm:projection}, \ref{thm:contraction}, \ref{thm:fixedpoint}, and \ref{thm:compactness}. Parameterwise non-collapse requires an additional anti-constant hypothesis introduced in Section~\ref{sec:parameterwise-noncollapse}.
\end{proof}

\section{The Projection Operator}

\begin{definition}[Metric Projection]\label{def:projection}
For $\Psi\in\Hilb$, define the metric projection onto $\I$ by
\[
\aleph(\Psi):=\arg\min_{\Phi\in\I}\norm{\Psi-\Phi}.
\]
\end{definition}

\begin{theorem}[Existence and Uniqueness of Projection]\label{thm:projection}
Under Hypothesis~\ref{hyp:H1}, the projection $\aleph(\Psi)$ exists and is unique for every $\Psi\in\Hilb$.
\end{theorem}

\begin{proof}
Existence follows because a minimizing sequence in the closed convex set $\I$ is Cauchy by the parallelogram law, hence converges in the complete space $\Hilb$ to a point of $\I$ that attains the infimum. Uniqueness follows from strict convexity of the norm induced by the Hilbert inner product: if two minimizers existed, their midpoint would produce the same minimum and the parallelogram identity would force the two minimizers to coincide.
\end{proof}

\begin{lemma}[Non-Expansiveness of the Projection]\label{lem:nonexp}
The map $\aleph:\Hilb\to\I$ is non-expansive:
\[
\norm{\aleph(\Psi_1)-\aleph(\Psi_2)}\le \norm{\Psi_1-\Psi_2}
\qquad
\forall \Psi_1,\Psi_2\in\Hilb.
\]
\end{lemma}

\begin{proof}
This is the standard Hilbert-space projection inequality. Let $\Phi_i=\aleph(\Psi_i)$. The projection characterization
\[
\inner{\Psi_i-\Phi_i}{\Phi-\Phi_i}\le 0
\qquad
\forall \Phi\in\I
\]
applied with $\Phi=\Phi_2$ and $\Phi=\Phi_1$ yields
\[
\norm{\Phi_1-\Phi_2}^2
\le
\inner{\Psi_1-\Psi_2}{\Phi_1-\Phi_2}
\le
\norm{\Psi_1-\Psi_2}\norm{\Phi_1-\Phi_2}.
\]
If $\Phi_1\neq \Phi_2$, divide by $\norm{\Phi_1-\Phi_2}$.
\end{proof}

\section{The Transition Operator}

\begin{definition}[Transition Operator]\label{def:transition}
For each $u\in\U$, define
\[
T_u(\Psi):=\aleph(K\Psi+\Gamma(u)).
\]
\end{definition}

\begin{theorem}[Contractivity]\label{thm:contraction}
Under Hypotheses~\ref{hyp:H1} and \ref{hyp:H2}, each $T_u$ is a strict contraction with Lipschitz constant at most $\norm{K}$:
\[
\norm{T_u(\Psi_1)-T_u(\Psi_2)}\le \norm{K}\,\norm{\Psi_1-\Psi_2}.
\]
\end{theorem}

\begin{proof}
By Lemma~\ref{lem:nonexp},
\begin{align*}
\norm{T_u(\Psi_1)-T_u(\Psi_2)}
&=
\norm{\aleph(K\Psi_1+\Gamma(u))-\aleph(K\Psi_2+\Gamma(u))} \\
&\le \norm{K(\Psi_1-\Psi_2)}
\le \norm{K}\,\norm{\Psi_1-\Psi_2}.
\end{align*}
Since $\norm{K}<1$, strict contraction follows.
\end{proof}

\section{Existence and Compactness of Attractors}

\begin{theorem}[Unique Fixed Point]\label{thm:fixedpoint}
Under Hypotheses~\ref{hyp:H1}--\ref{hyp:H3}, for each $u\in\U$ there exists a unique $f_u^\ast\in\Hilb$ such that
\[
T_u(f_u^\ast)=f_u^\ast.
\]
Moreover, for every $\Psi_0\in\Hilb$,
\[
\lim_{n\to\infty}T_u^n(\Psi_0)=f_u^\ast.
\]
\end{theorem}

\begin{proof}
By Theorem~\ref{thm:contraction}, each $T_u$ is a strict contraction on the complete metric space $(\Hilb,\norm{\cdot})$. The Banach fixed-point theorem yields existence, uniqueness, and convergence of iterates.
\end{proof}

\begin{definition}[Attractor Family]\label{def:attractor}
Define the attractor family
\[
\A^\ast:=\{f_u^\ast:u\in\U\}.
\]
\end{definition}

\begin{theorem}[Compactness of the Attractor Family]\label{thm:compactness}
Under Hypotheses~\ref{hyp:H1}--\ref{hyp:H4}, the set $\A^\ast$ is compact in $\Hilb$.
\end{theorem}

\begin{proof}
Let $F:\U\to\Hilb$ be given by $F(u)=f_u^\ast$. The fixed points depend continuously on $u$ because the contraction constant is uniform and $\Gamma$ is uniformly continuous on the compact parameter space. Since $\U$ is compact and $\A^\ast=F(\U)$, the image is compact.
\end{proof}

\section{Parameterwise Non-Collapsibility}\label{sec:parameterwise-noncollapse}

\begin{definition}[Constant Subspace]
Whenever the ambient realization distinguishes constant elements, let $\mathcal{N}\subset\Hilb$ denote the closed subspace of constants.
\end{definition}

\begin{definition}[Parameterwise Non-Collapsible Family]
A family $\{T_u\}_{u\in\U}$ is parameterwise non-collapsible if
\[
f_u^\ast\notin\mathcal{N}
\qquad
\forall u\in\U.
\]
\end{definition}

\begin{hypothesis}[H5: No Constant Fixed Points]\label{hyp:H5}
For every parameter $u\in\U$ and every constant element $c\in\mathcal{N}$,
\[
T_u(c)\neq c.
\]
\end{hypothesis}

\begin{theorem}[Structural Non-Collapsibility]\label{thm:noncollapse}
Under Hypotheses~\ref{hyp:H1}--\ref{hyp:H5}, the family is parameterwise non-collapsible:
\[
f_u^\ast\notin\mathcal{N}
\qquad
\forall u\in\U.
\]
\end{theorem}

\begin{proof}
Fix $u\in\U$. Suppose by contradiction that $f_u^\ast\in\mathcal{N}$. Write $f_u^\ast=c$ with $c\in\mathcal{N}$. Since $f_u^\ast$ is a fixed point,
\[
T_u(c)=T_u(f_u^\ast)=f_u^\ast=c,
\]
contradicting Hypothesis~\ref{hyp:H5}. Therefore $f_u^\ast\notin\mathcal{N}$.
\end{proof}

\begin{remark}[What H5 does not say]
Hypothesis~\ref{hyp:H5} is parameterwise. It does not merely say that \emph{some} parameter breaks constant collapse. It excludes constant fixed points for \emph{every} parameter. This is exactly the quantifier strength needed for Theorem~\ref{thm:noncollapse}.
\end{remark}

\section{Summary of Base Results}

\begin{enumerate}[leftmargin=*]
\item H1--H4 give a complete theorem-tight base for projection existence, strict contraction, unique fixed points, and compact attractor families.
\item H5 is a separate anti-constant assumption block used only for parameterwise non-collapse.
\item No empirical, phenomenological, biological, or runtime claim is inferred at this level.
\item Later sections may add conditional identity, rigidity, and criterion theorems, but only by introducing their own explicit assumption blocks.
\end{enumerate}

\part{Typed Computable Model and State Spectral Gap}
\section{Typed Computable Model: Embedded Constant-Disk Witness}

This section repairs the type leaks of the previous computable model. The purpose is narrow: provide a fully typed witness for the transition grammar and the state spectral gap. It does \emph{not} by itself certify parameterwise non-collapse; that burden is handled separately below.

\subsection{Golden Mean Parameter Space}

\begin{definition}[Golden Mean Shift Parameter Space]
Let $\Sigma=\{0,1\}$ and
\[
A=
\begin{pmatrix}
1 & 1\\
1 & 0
\end{pmatrix}.
\]
The Golden Mean Shift $S_A\subset \Sigma^{\mathbb{N}}$ is the subshift of finite type consisting of all binary sequences that avoid the block $11$.
\end{definition}

\begin{definition}[Parameter Encoding]
Define
\[
u(x)=\sum_{i=1}^{\infty} x_i\phi^{-i},
\qquad
\phi=\frac{1+\sqrt{5}}{2}.
\]
The image $\U:=u(S_A)\subset [0,1]$ is compact.
\end{definition}

\subsection{Typed Banach Witness}

\begin{definition}[Ambient Space and Constant Embedding]
Let
\[
X=[0,1],
\qquad
E=\mathbb{R}^2,
\qquad
\mathcal{B}=C(X,E)
\]
with norm
\[
\norm{f}_\infty:=\sup_{s\in X}\norm{f(s)}_2.
\]
Define the constant embedding
\[
j:E\to \mathcal{B},
\qquad
j(v)(s)=v.
\]
\end{definition}

\begin{definition}[Invariant Constant Disk]
Let
\[
D:=\{v\in E:\norm{v}_2\le 1\},
\qquad
\mathcal{I}_{\mathrm{c}}:=j(D)\subset \mathcal{B}.
\]
\end{definition}

\begin{definition}[Averaging Operator]
Define
\[
A:\mathcal{B}\to E,
\qquad
A(f):=\int_0^1 f(\tau)\,d\tau.
\]
\end{definition}

\begin{definition}[Disk Projection]
Define the Euclidean projection
\[
P_E:E\to D,
\qquad
P_E(v)=
\begin{cases}
v, & \norm{v}_2\le 1,\\[0.4em]
\dfrac{v}{\norm{v}_2}, & \norm{v}_2>1.
\end{cases}
\]
\end{definition}

\begin{definition}[Typed Projection onto the Constant Disk]
Define
\[
\aleph_{\mathcal B}:\mathcal{B}\to \mathcal{I}_{\mathrm{c}},
\qquad
\aleph_{\mathcal B}(f):=j(P_E(A(f))).
\]
\end{definition}

\begin{definition}[Constant-Valued Contraction Operator]
Define
\[
\widetilde K:\mathcal{B}\to\mathcal{B},
\qquad
(\widetilde K f)(s):=\frac{1}{2}\int_0^1 f(\tau)\,d\tau.
\]
\end{definition}

\begin{definition}[Lifted Input Coupling]
Let $\Gamma:\U\to E$ be
\[
\Gamma(u)=(u,0).
\]
Define its constant lift
\[
\Gamma_{\mathcal B}:\U\to \mathcal{B},
\qquad
\Gamma_{\mathcal B}(u):=j(\Gamma(u)).
\]
\end{definition}

\begin{definition}[Typed Transition Operator]
For each $u\in\U$, define
\[
T_u:\mathcal{B}\to \mathcal{I}_{\mathrm{c}},
\qquad
T_u(f):=\aleph_{\mathcal B}\bigl(\widetilde K f+\Gamma_{\mathcal B}(u)\bigr).
\]
\end{definition}

\begin{remark}[Typing discipline]
No unembedded element of $E$ is treated as an element of $\mathcal{B}$. The embedding $j$ is always written explicitly. This removes the earlier type leak in which a disk in $E$ was silently treated as a subset of a function space.
\end{remark}

\subsection{Basic Operator Properties}

\begin{lemma}[The Embedding is Isometric]
For every $v\in E$,
\[
\norm{j(v)}_\infty=\norm{v}_2.
\]
\end{lemma}

\begin{proof}
By definition,
\[
\norm{j(v)}_\infty=\sup_{s\in X}\norm{j(v)(s)}_2
=\sup_{s\in X}\norm{v}_2
=\norm{v}_2.
\]
\end{proof}

\begin{lemma}[The Disk Projection is Non-Expansive]\label{lem:disk-projection}
For all $v,w\in E$,
\[
\norm{P_E(v)-P_E(w)}_2\le \norm{v-w}_2.
\]
\end{lemma}

\begin{proof}
$P_E$ is the metric projection onto the closed convex set $D\subset \mathbb{R}^2$, so the standard Hilbert-space non-expansiveness property applies.
\end{proof}

\begin{lemma}[The Typed Projection is Non-Expansive]
For all $f,g\in\mathcal{B}$,
\[
\norm{\aleph_{\mathcal B}(f)-\aleph_{\mathcal B}(g)}_\infty\le \norm{f-g}_\infty.
\]
\end{lemma}

\begin{proof}
Using the isometry of $j$, the non-expansiveness of $P_E$, and Jensen's inequality for the averaging operator,
\begin{align*}
\norm{\aleph_{\mathcal B}(f)-\aleph_{\mathcal B}(g)}_\infty
&=
\norm{P_E(A(f))-P_E(A(g))}_2 \\
&\le \norm{A(f)-A(g)}_2 \\
&=
\left\|\int_0^1 (f(\tau)-g(\tau))\,d\tau\right\|_2 \\
&\le \int_0^1 \norm{f(\tau)-g(\tau)}_2\,d\tau
\le \norm{f-g}_\infty.
\end{align*}
\end{proof}

\begin{lemma}[Norm of the Lifted Contraction]
\[
\norm{\widetilde K}\le \frac{1}{2}.
\]
\end{lemma}

\begin{proof}
For any $f\in\mathcal{B}$,
\[
\norm{\widetilde Kf}_\infty
=
\left\|\frac{1}{2}\int_0^1 f(\tau)\,d\tau\right\|_2
\le \frac{1}{2}\int_0^1 \norm{f(\tau)}_2\,d\tau
\le \frac{1}{2}\norm{f}_\infty.
\]
\end{proof}

\begin{theorem}[Typed Contraction on the Computable Witness]
For each $u\in\U$, the map $T_u:\mathcal{B}\to\mathcal{I}_{\mathrm c}$ is a strict contraction:
\[
\norm{T_u(f)-T_u(g)}_\infty\le \frac{1}{2}\norm{f-g}_\infty.
\]
\end{theorem}

\begin{proof}
Combine non-expansiveness of $\aleph_{\mathcal B}$ with the bound on $\widetilde K$:
\[
\norm{T_u(f)-T_u(g)}_\infty
\le
\norm{\widetilde K f-\widetilde K g}_\infty
\le
\frac{1}{2}\norm{f-g}_\infty.
\]
\end{proof}

\begin{corollary}[Unique Fixed Point in the Typed Witness]
For each $u\in\U$, there exists a unique $f_u^\ast\in\mathcal{I}_{\mathrm c}$ such that
\[
T_u(f_u^\ast)=f_u^\ast.
\]
\end{corollary}

\begin{proof}
The Banach fixed-point theorem applies on the complete metric space $(\mathcal{B},\norm{\cdot}_\infty)$.
\end{proof}

\begin{proposition}[Typed Non-Vacuity of the Computable Model]\label{prop:Q3verify}
The embedded constant-disk Golden Mean model is a fully typed, non-vacuous witness of the transition grammar:
\begin{enumerate}[leftmargin=*]
\item every operator has explicit domain and codomain;
\item the transition map $T_u$ is well-defined on $\mathcal{B}$;
\item the family admits unique fixed points and a compact parameter image;
\item the state contraction bound is explicit and computable.
\end{enumerate}
\end{proposition}

\begin{proof}
Items (1)--(3) follow from the definitions above and the contraction theorem. Compactness of the parameter image follows because $\U$ is compact and $u\mapsto f_u^\ast$ is continuous under a uniform contraction constant.
\end{proof}

\begin{remark}[Why this witness does not certify H5]
The image of $T_u$ lies in the constant subspace $\mathcal{I}_{\mathrm c}$. Therefore its fixed point is constant by construction. This witness is correct and useful for typing and state spectral analysis, but it cannot certify the parameterwise anti-constant hypothesis H5 from Section~\ref{sec:parameterwise-noncollapse}.
\end{remark}

\subsection{Separate Anti-Collapse Witness}

\begin{definition}[Affine Non-Constant Witness Space]
Let
\[
\Hilb_{\mathrm{ac}}:=L^2([0,1]),
\qquad
\mathcal{N}:=\{c\mathbf{1}:c\in\mathbb{R}\},
\qquad
h(s):=\sin(2\pi s).
\]
Choose $r$ with $0<r<\norm{h}_{L^2}$ and define the closed convex affine set
\[
\I_{\mathrm{ac}}:=h+\{g\in\Hilb_{\mathrm{ac}}:\norm{g}_{L^2}\le r\}.
\]
\end{definition}

\begin{lemma}[The Anti-Collapse Witness Set Avoids Constants]
\[
\I_{\mathrm{ac}}\cap \mathcal{N}=\varnothing.
\]
\end{lemma}

\begin{proof}
The distance from $h$ to the constant subspace $\mathcal{N}$ equals $\norm{h}_{L^2}$ because $h$ has mean zero and is orthogonal to constants. Since $r<\norm{h}_{L^2}$, the radius-$r$ ball around $h$ does not intersect $\mathcal{N}$.
\end{proof}

\begin{definition}[Anti-Collapse Witness Dynamics]
Let $\U=[0,1]$, choose $\beta\in(0,1/2)$, and define
\[
K_{\mathrm{ac}}f:=\beta f,
\qquad
\Gamma_{\mathrm{ac}}(u):=\frac{u}{4}h.
\]
Let $P_{\mathrm{ac}}$ denote the metric projection onto $\I_{\mathrm{ac}}$, and define
\[
T_u^{\mathrm{ac}}(f):=P_{\mathrm{ac}}(K_{\mathrm{ac}}f+\Gamma_{\mathrm{ac}}(u)).
\]
\end{definition}

\begin{proposition}[Separate Witness for H5]
The family $\{T_u^{\mathrm{ac}}\}_{u\in\U}$ satisfies Hypotheses~\ref{hyp:H1}--\ref{hyp:H5} of Section~1.
\end{proposition}

\begin{proof}
\begin{enumerate}[leftmargin=*]
\item $\Hilb_{\mathrm{ac}}$ is a Hilbert space and $\I_{\mathrm{ac}}$ is non-empty, closed, and convex.
\item $\norm{K_{\mathrm{ac}}}=\beta<1$.
\item Completeness is inherited from $L^2([0,1])$.
\item $\U=[0,1]$ is compact and $\Gamma_{\mathrm{ac}}$ is continuous.
\item For every constant $c\in\mathcal{N}$ and every $u\in\U$, the value $T_u^{\mathrm{ac}}(c)$ belongs to $\I_{\mathrm{ac}}$, while $\I_{\mathrm{ac}}\cap\mathcal{N}=\varnothing$. Hence $T_u^{\mathrm{ac}}(c)\neq c$.
\end{enumerate}
\end{proof}

\begin{remark}
BaseCore therefore distinguishes two roles cleanly: the Golden Mean witness certifies typed operator correctness and state contraction; the affine anti-collapse witness certifies the stronger H5 burden needed for parameterwise non-collapse.
\end{remark}

\section{State Spectral Gap and Parameter-Family Persistence}

\begin{theorem}[State Spectral Gap]\label{thm:spectrum}
Fix $u\in\U$ in the typed computable model. Assume $T_u$ is Fr\'echet differentiable at its fixed point $f_u^\ast$. Then
\[
\rho\!\left(DT_u(f_u^\ast)\right)
\le
\norm{DT_u(f_u^\ast)}
\le
\norm{\widetilde K}
\le
\frac{1}{2}.
\]
Consequently the state spectral gap satisfies
\[
\Delta_{\mathrm{state}}\ge 1-\norm{\widetilde K}\ge \frac{1}{2}.
\]
\end{theorem}

\begin{proof}
At every point where the projection $P_E$ is Fr\'echet differentiable, its derivative has operator norm at most $1$ because $P_E$ is non-expansive. Since $\aleph_{\mathcal B}=j\circ P_E\circ A$,
\[
DT_u(f_u^\ast)=
D\aleph_{\mathcal B}\!\left(\widetilde K f_u^\ast+\Gamma_{\mathcal B}(u)\right)\circ \widetilde K,
\]
and therefore
\[
\norm{DT_u(f_u^\ast)}\le \norm{D\aleph_{\mathcal B}}\norm{\widetilde K}\le \frac{1}{2}.
\]
The spectral-radius bound follows from $\rho(L)\le \norm{L}$ for bounded linear operators.
\end{proof}

\begin{theorem}[Relaxation-Time Bound]\label{thm:relaxation}
Under the hypotheses of Theorem~\ref{thm:spectrum}, any discrete-to-continuous relaxation convention of the form
\[
\tau_{\mathrm{state}}=-\frac{1}{\ln \rho}
\]
with $\rho=\rho(DT_u(f_u^\ast))$ gives
\[
\tau_{\mathrm{state}}\le \frac{1}{\ln 2}.
\]
\end{theorem}

\begin{proof}
Theorem~\ref{thm:spectrum} gives $\rho\le 1/2$. Since $-\ln(\rho)\ge \ln 2$, the bound follows.
\end{proof}

\begin{proposition}[Parameter-Family Persistence is not State Spectrum]
Let $F:\U\to\mathcal{B}$ be the fixed-point family $F(u)=f_u^\ast$. Any persistence of parameter labels or family coordinates belongs to the variation of $F$ over $\U$, not to the state-transition spectrum of $DT_u(f_u^\ast)$ for fixed $u$.
\end{proposition}

\begin{proof}
The state-transition spectrum is defined for the linearization of $T_u$ at fixed parameter $u$. Family-level transport instead concerns the map $F$ between parameter space and fixed-point set. These are different operators on different domains. Therefore a family-persistence mode cannot be used to introduce a unit eigenvalue into the contraction spectrum of $DT_u(f_u^\ast)$.
\end{proof}

\begin{remark}[What has been removed]
The previous text incorrectly mixed strict contraction with a state-spectrum eigenvalue $\lambda=1$. BaseCore no longer does that. Any neutral or persistent mode must now be interpreted as a parameter-family bookkeeping mode, not as an eigenmode of the contractive state Jacobian.
\end{remark}

\part{Identity, Rigidity, and Conditional Non-Simulability}
\section{Projective Systems and Inverse-Limit Identity}

\subsection{Projective Setup}

\begin{definition}[Projective System of State Spaces]
Let $\{S_t\}_{t\in\mathbb{N}}$ be a family of non-empty state spaces together with projection maps
\[
\pi_{t+1\to t}:S_{t+1}\to S_t
\]
satisfying the compatibility rule
\[
\pi_{t+1\to t}\circ \pi_{t+2\to t+1}=\pi_{t+2\to t}
\qquad
\forall t\in\mathbb{N}.
\]
\end{definition}

\begin{definition}[Identity as Inverse Limit]
The identity object associated with the projective system is
\[
\Id:=\varprojlim S_t
=
\left\{
(x_t)_{t\in\mathbb{N}}\in \prod_{t\in\mathbb{N}} S_t
\;\middle|\;
\pi_{t+1\to t}(x_{t+1})=x_t\ \forall t
\right\}.
\]
\end{definition}

\begin{remark}[Identity neutrality]
In BaseCore, $\Id$ is a structural inverse-limit object only. It carries no automatic phenomenological, biological, or metaphysical interpretation.
\end{remark}

\begin{hypothesis}[Topological Regularity]\label{hyp:topo}
Each $S_t$ is compact Hausdorff and each $\pi_{t+1\to t}$ is continuous and surjective.
\end{hypothesis}

\begin{hypothesis}[NFD: Non-Finite Determination]\label{hyp:nfd}
For every $t\in\mathbb{N}$, the canonical projection
\[
p_t:\Id\to S_t,
\qquad
p_t((x_k)_k)=x_t
\]
is not injective. Equivalently,
\[
\forall t\ \exists x,y\in\Id
\quad
x\neq y
\quad\text{and}\quad
p_t(x)=p_t(y).
\]
\end{hypothesis}

\begin{remark}[Persistent future branching as an operational form of NFD]
A sufficient operational strengthening of Hypothesis~\ref{hyp:nfd} is: for every $t$ there exists $k>t$ and $x_t\in S_t$ such that the fiber of $\pi_{k\to t}$ over $x_t$ contains at least two globally extendable compatible histories. BaseCore keeps the weaker NFD statement as the primary assumption because it is exactly what the non-locality theorem needs.
\end{remark}

\subsection{Non-Locality of Identity}

\begin{theorem}[Non-Locality of Identity under NFD]\label{prop:nonlocality}
Assume Hypotheses~\ref{hyp:topo} and \ref{hyp:nfd}. Then no canonical finite-time observation determines identity. More precisely, for every $t$ the canonical projection
\[
p_t:\Id\to S_t
\]
is not injective, so identity is not recoverable from any finite slice through the projective observation maps.
\end{theorem}

\begin{proof}
This is exactly Hypothesis~\ref{hyp:nfd}. The theorem restates the structural consequence in claim-boundary language: finite slices can participate in the inverse-limit object, but no single canonical slice determines it.
\end{proof}

\begin{remark}[What is \emph{not} proved here]
BaseCore does not prove that no arbitrary set-theoretic injection $\Id\to S_t$ can exist in every imaginable inverse system. It proves the narrower and correct claim that the canonical projective observations $p_t$ are not injective under Hypothesis~\ref{hyp:nfd}.
\end{remark}

\begin{proposition}[No Canonical Progressive Recovery under NFD]\label{thm:noprogressive}
Assume Hypothesis~\ref{hyp:nfd}. No finite canonical prefix map determines the inverse-limit identity object.
\end{proposition}

\begin{proof}
If some finite canonical prefix determined identity, then in particular one of the canonical coordinate maps would become injective after finite truncation, contradicting Hypothesis~\ref{hyp:nfd}.
\end{proof}

\section{Rigidity under Admissible Perturbations}

\subsection{Metric Projective Assumptions}

\begin{assumption}[RIG-1: Metric Projective System]\label{ass:rig1}
Each $S_t$ is equipped with a complete metric $d_t$.
\end{assumption}

\begin{assumption}[RIG-2: Uniform Lipschitz Projections]\label{ass:rig2}
Each projection $\pi_{t+1\to t}:S_{t+1}\to S_t$ is Lipschitz with constant $L_t$, and the weighted control sequence used below is finite.
\end{assumption}

\begin{assumption}[RIG-3: Stable Lifting]\label{ass:rig3}
For each admissible perturbed projection family there exist compatible approximate lifting maps with constants bounded uniformly by a finite lifting constant $L_{\mathrm{lift}}$.
\end{assumption}

\begin{assumption}[RIG-4: Typed Perturbation Model]\label{ass:rig4}
Perturbations are represented either by maps between common ambient Banach embeddings or by map-space distances
\[
d_{\mathrm{map}}^{(t)}(\pi_{t+1\to t},\widetilde\pi_{t+1\to t})\le \eps_t.
\]
Literal expressions of the form $\widetilde\pi=\pi+\eps$ are used only when a common ambient linear model has been specified.
\end{assumption}

\begin{assumption}[RIG-5: Summable Deformation]\label{ass:rig5}
There exists a positive weight sequence $\{w_t\}_{t\ge 0}$ such that
\[
\sum_{t=0}^\infty w_t\eps_t<\infty.
\]
\end{assumption}

\begin{definition}[Weighted Deformation Size]
The weighted deformation size of a perturbed projective family is
\[
\|\eps_\bullet\|_w:=\sum_{t=0}^{\infty} w_t\eps_t.
\]
\end{definition}

\begin{definition}[Weighted Product Metric]
For compatible sequences $x=(x_t)$ and $y=(y_t)$ define
\[
\dw(x,y):=\sup_{t\ge 0} w_t\, d_t(x_t,y_t).
\]
\end{definition}

\begin{theorem}[Conditional Stability of Inverse-Limit Identity]\label{thm:rigidity}
Assume RIG-1 through RIG-5. Then the perturbed inverse limit
\[
\widetilde{\Id}:=\varprojlim \widetilde S_t
\]
exists, and there is a constant $C_{\mathrm{rig}}>0$ depending only on the projection, lifting, and weight controls such that
\[
d_H(\Id,\widetilde{\Id})\le C_{\mathrm{rig}}\|\eps_\bullet\|_w
\]
in the weighted product metric.
\end{theorem}

\begin{proof}[Proof sketch]
The lifting assumption gives compatible approximate lifts at each level. Weighted Lipschitz control propagates local deformation bounds through the projective tower. Summability ensures that the cumulative compatibility defect remains finite, so diagonal compactness arguments produce a perturbed inverse-limit object. The same bounds control the Hausdorff distance between the original and perturbed compatible-sequence sets.
\end{proof}

\begin{corollary}[Conditional Isomorphism Requires Extra Structure]
The metric estimate of Theorem~\ref{thm:rigidity} does not by itself imply that $\Id$ and $\widetilde{\Id}$ are isomorphic. Isomorphism or shape-equivalence requires additional bi-lifting, coherence, or ambient-embedding assumptions.
\end{corollary}

\begin{proof}
Metric closeness and categorical isomorphism are different conclusions. A small Hausdorff deformation need not preserve global inverse-limit type without extra hypotheses guaranteeing two-sided coherent transport.
\end{proof}

\begin{definition}[Ontological Mass]\label{def:ontmass}
Let $E_w(\delta)$ denote the weighted energy of an admissible perturbation family $\delta$. Define
\[
\MO:=\inf\left\{
\frac{E_w(\delta)}{d_H(\Id,\widetilde{\Id}_\delta)}
\;\middle|\;
\delta \text{ admissible},\ d_H(\Id,\widetilde{\Id}_\delta)>0
\right\},
\]
with the convention $\MO=+\infty$ if no admissible finite-energy perturbation yields positive deformation of the identity object.
\end{definition}

\begin{remark}[Interpretation discipline]
$\MO$ is a deformation-resistance scalar, not a metaphysical substance. BaseCore uses it only as a conditional invariant attached to admissible perturbation models.
\end{remark}

\section{Conditional Non-Simulability}

\begin{assumption}[NS-1: CCR Target]\label{ass:ns1}
A CCR target is an inverse-limit identity object satisfying:
\begin{enumerate}[leftmargin=*]
\item non-finite determination (Hypothesis~\ref{hyp:nfd});
\item infinite deformation resistance $\MO=+\infty$;
\item absence of a finite sufficient statistic for faithful identity recovery;
\item absence of a finite-dimensional faithful embedding preserving projective compatibility.
\end{enumerate}
\end{assumption}

\begin{assumption}[NS-2: Finite Simulator Class]\label{ass:ns2}
A simulator $C$ belongs to $\mathbf{Comp}_{\mathrm{finite}}$ if it has at least one of the following boundedness features:
\begin{enumerate}[leftmargin=*]
\item finite-dimensional latent state or finite context;
\item bounded memory horizon or finite sufficient statistic;
\item finite perturbation-energy budget;
\item finite-rank projective representation.
\end{enumerate}
\end{assumption}

\begin{assumption}[NS-3: Faithful Simulation Requirement]\label{ass:ns3}
A simulation is faithful only if it preserves, at the stated tolerance level, all of:
\begin{enumerate}[leftmargin=*]
\item projective compatibility;
\item non-finite determination;
\item deformation-resistance class;
\item relevant observables in the weighted metric;
\item admissible intervention responses.
\end{enumerate}
\end{assumption}

\begin{theorem}[Conditional Non-Simulability of CCR by Finite Simulators]\label{thm:non-simulability}
Assume NS-1 through NS-3. Then no simulator in $\mathbf{Comp}_{\mathrm{finite}}$ faithfully realizes a CCR target.
\end{theorem}

\begin{proof}
A finite simulator in the sense of NS-2 necessarily compresses the target into finite latent structure, finite memory horizon, finite-rank projective representation, or a finite sufficient statistic. Any such compression destroys at least one component of NS-1: either NFD collapses, the faithful projective structure is lost, or the deformation-resistance class is reduced from $\MO=+\infty$ to a finite value. Therefore faithful realization, as defined by NS-3, fails.
\end{proof}

\begin{theorem}[Approximation Barrier]
Assume NS-1 through NS-3. Finite simulators may approximate bounded finite projections of a CCR target over finite horizons, but they cannot preserve the full inverse-limit object faithfully as the horizon grows without bound unless resources also grow without bound.
\end{theorem}

\begin{proof}
Finite-horizon projections are finite objects and may be approximated by finite representations. But a faithful approximation of the full inverse-limit target must preserve NFD and the infinite deformation-resistance class. That cannot be achieved with bounded representational resources by Theorem~\ref{thm:non-simulability}.
\end{proof}

\begin{remark}[What has been removed]
BaseCore no longer claims that finite simulators ``cannot approximate at all.'' The correct statement is conditional and finer-grained: finite-horizon approximation may exist, but faithful full inverse-limit realization of CCR targets is blocked under NS-1 through NS-3.
\end{remark}

\section{Section Summary}

\begin{enumerate}[leftmargin=*]
\item Identity is an inverse-limit object, not a canonical finite slice.
\item Non-locality now depends explicitly on the NFD assumption.
\item Rigidity is stated as metric stability under explicit perturbation-model assumptions, not as automatic isomorphism.
\item Non-simulability is conditional on a defined CCR target, a defined finite simulator class, and a defined notion of faithful simulation.
\item BaseCore therefore preserves theorem strength where assumptions support it and weakens only those statements that were previously over-absolute.
\end{enumerate}

% Absorbed from paper2/main.tex
\part{Regime Structure and Continuity Constraints}
\section{Formal Setup}

\subsection{Review of Identity Framework}

We assume the identity framework of \cite{paper1}. Let $(S_t, \pi_{t+1 \to t})_{t \in \mathbb{N}}$ be a projective system of observable state spaces. The identity object is the inverse limit:
\begin{equation}
    \mathcal{I} := \varprojlim S_t = \left\{ (x_t)_t \in \prod_t S_t \;\Big|\; \pi_{t+1 \to t}(x_{t+1}) = x_t \;\; \forall t \right\}.
\end{equation}

The ontological mass $M_\Omega$ quantifies resistance to identity deformation. We recall that:
\begin{itemize}[leftmargin=*]
    \item $M_\Omega = 0$: null persistence (no stable identity);
    \item $0 < M_\Omega < \infty$: deformable persistence (finite rigidity);
    \item $M_\Omega = +\infty$: closed--rigid (CCR, absolute rigidity).
\end{itemize}

\subsection{Abstract Phenomenological Space}

\begin{definition}[Phenomenological Space]\label{def:phenspace}
A phenomenological space is a nonempty set $\mathcal{E}$ whose elements are called \textit{phenomenological configurations}. No topology, metric, or algebraic structure is assumed a priori.
\end{definition}

\begin{remark}
We do not claim that $\mathcal{E}$ represents ``conscious states'' or ``qualia''. It is an abstract mathematical space whose interpretation is not specified by the formalism. The space $\mathcal{E}$ is introduced purely as the codomain of a structural assignment; its ontological status is deliberately left open.
\end{remark}

\subsubsection{Why a Separate Space?}

One might ask: why not use $\mathcal{I}$ itself as the phenomenological space? The answer is structural:

\begin{proposition}[Non-Triviality of Phenomenological Codomain]
If phenomenological assignments were required to land in $\mathcal{I}$ itself, then any compatible assignment would be the identity map, eliminating all non-trivial phenomenological structure.
\end{proposition}

\begin{proof}
Let $\Phi: \mathcal{I} \to \mathcal{I}$ be compatible with projections. By the universal property of inverse limits, $\Phi$ must commute with all projection maps. The only such endomorphism on $\mathcal{I}$ is the identity (up to isomorphism). Hence no non-trivial phenomenological stratification is possible.
\end{proof}

Thus, a separate space $\mathcal{E}$ is \textit{structurally necessary} for non-trivial phenomenological content.

\subsection{Phenomenological Assignment}

\begin{definition}[Phenomenological Assignment]\label{def:phenassign}
A phenomenological assignment is a family of partial maps:
\begin{equation}
    \phi_t: S_t \to \mathcal{E}
\end{equation}
indexed by $t \in \mathbb{N}$.
\end{definition}

\begin{definition}[Induced Limit Assignment]\label{def:limitassign}
Given a phenomenological assignment $\{\phi_t\}$, the induced limit assignment is the partial map:
\begin{equation}
    \Phi: \mathcal{I} \to \mathcal{E}, \quad \Phi((x_t)_t) := \lim_{t \to \infty} \phi_t(x_t),
\end{equation}
whenever the limit exists in an appropriate sense (to be defined by compatibility axioms).
\end{definition}

\subsubsection{Uniqueness of the Limit Assignment}

\begin{proposition}[Uniqueness of Compatible $\Phi$]\label{prop:unique-phi}
Given a compatible family $\{\phi_t\}$ satisfying E1, the induced limit assignment $\Phi$ is unique.
\end{proposition}

\begin{proof}
By E1, $\phi_{t+1}(x_{t+1}) = \phi_t(x_t)$ for all $t$. Hence the sequence $\{\phi_t(x_t)\}$ is constant. The limit exists trivially and equals this constant value. Uniqueness follows from the constancy of the sequence.
\end{proof}

\subsection{Axioms for Phenomenological Coherence}

We introduce minimal axioms under which a phenomenological assignment is \textit{structurally admissible}. We will then prove that these axioms are not arbitrary choices but structurally forced conditions.

\begin{axiom}[E0: Non-Triviality]\label{ax:e0}
The phenomenological space satisfies $|\mathcal{E}| \geq 2$, i.e., there exist at least two distinct phenomenological configurations.
\end{axiom}

\begin{axiom}[E1: Projective Compatibility]\label{ax:e1}
For every compatible identity trajectory $(x_t)_t \in \mathcal{I}$, the induced phenomenological sequence satisfies:
\begin{equation}
    \phi_{t+1}(x_{t+1}) = \phi_t(x_t) \quad \forall t.
\end{equation}
\end{axiom}

\begin{axiom}[E2: Non-Locality]\label{ax:e2}
There exists no function $f: S_t \to \mathcal{E}$ (for any fixed $t$) such that for all trajectories $(x_t)_t \in \mathcal{I}$:
\begin{equation}
    \Phi((x_t)_t) = f(x_t).
\end{equation}
\end{axiom}

\subsection{Inevitability of the Axioms}

We now prove that axioms E0--E2 are \textit{not design choices} but structurally necessary conditions for any meaningful notion of phenomenological assignment.

\subsubsection{Inevitability of E0}

\begin{theorem}[E0 is Necessary]\label{thm:e0-necessary}
If $|\mathcal{E}| = 1$, then every phenomenological assignment is constant and no phenomenological classification is possible.
\end{theorem}

\begin{proof}
If $\mathcal{E} = \{e_0\}$, then $\Phi(x) = e_0$ for all $x \in \mathcal{I}$. Hence:
\begin{enumerate}[leftmargin=*]
    \item No phenomenological distinction between identity trajectories exists;
    \item The Fragmentation and Forced Continuity Theorems become vacuous;
    \item The downstream null-regime closure developed in \textit{Null-Regime Instability} would have no content.
\end{enumerate}
Thus, $|\mathcal{E}| \geq 2$ is necessary for the framework to have non-trivial content.
\end{proof}

\subsubsection{Inevitability of E1}

\begin{theorem}[E1 is Necessary]\label{thm:e1-necessary}
If E1 fails, then phenomenological assignments are incompatible with the inverse-limit structure of identity, and no coherent $\Phi: \mathcal{I} \to \mathcal{E}$ can be defined.
\end{theorem}

\begin{proof}
Suppose E1 fails for some trajectory $(x_t)_t \in \mathcal{I}$. Then there exists $t_0$ such that:
\begin{equation}
    \phi_{t_0+1}(x_{t_0+1}) \neq \phi_{t_0}(x_{t_0}).
\end{equation}
But the identity trajectory satisfies $\pi_{t_0+1 \to t_0}(x_{t_0+1}) = x_{t_0}$. Hence the phenomenological assignment fails to respect the projective structure that defines identity.

This means $\Phi$ cannot be factored through the inverse limit:
\begin{equation}
    \Phi \neq \varprojlim \phi_t.
\end{equation}
Hence no well-defined limit assignment exists. Without E1, the phenomenological assignment is structurally incompatible with identity.
\end{proof}

\subsubsection{Inevitability of E2}

\begin{theorem}[E2 is Necessary]\label{thm:e2-necessary}
If E2 fails, then phenomenology is a local property of states, contradicting the global nature of identity as an inverse limit.
\end{theorem}

\begin{proof}
Suppose E2 fails. Then there exists $t_0$ and $f: S_{t_0} \to \mathcal{E}$ such that:
\begin{equation}
    \Phi((x_t)_t) = f(x_{t_0}) \quad \forall (x_t)_t \in \mathcal{I}.
\end{equation}
This implies:
\begin{enumerate}[leftmargin=*]
    \item Phenomenology is determined by a finite temporal slice;
    \item The full trajectory $(x_t)_{t > t_0}$ is irrelevant to $\Phi$;
    \item Identity as an inverse limit has no phenomenological content beyond $t_0$.
\end{enumerate}

But identity is defined as a \textit{global} object over infinite time. If phenomenology is local, it cannot reflect the global structure of identity. This contradicts the foundational premise that phenomenological assignments should be compatible with identity structure.

Hence E2 is necessary for phenomenology to be structurally aligned with inverse-limit identity.
\end{proof}

\subsection{Minimality Theorem}

\begin{theorem}[Axiomatic Minimality]\label{thm:minimality}
The axiom set $\{E0, E1, E2\}$ is minimal: no proper subset is sufficient to derive the main results of this paper.
\end{theorem}

\begin{proof}
We show that removing any axiom destroys the framework:

\textbf{Without E0:} All assignments are trivially constant. Fragmentation Theorem (4.1) is vacuous.

\textbf{Without E1:} No coherent $\Phi$ exists. The entire framework collapses.

\textbf{Without E2:} Phenomenology becomes local. The Forced Continuity Theorem (4.2) fails because local functions can be arbitrarily perturbed.

Hence $\{E0, E1, E2\}$ is the minimal viable axiom set.
\end{proof}

\subsection{Rejected Alternatives}

We briefly discuss axiom candidates that were considered and rejected.

\subsubsection{Rejected: Continuity Axiom}
One might require $\Phi$ to be continuous. This is \textit{not} an axiom because:
\begin{itemize}[leftmargin=*]
    \item It requires pre-existing topology on $\mathcal{E}$;
    \item The Forced Continuity Theorem \textit{derives} continuity from CCR, not as assumption.
\end{itemize}

\subsubsection{Rejected: Surjectivity Axiom}
One might require $\Phi$ to be surjective. This is \textit{not} an axiom because:
\begin{itemize}[leftmargin=*]
    \item It would exclude null assignments, preventing the downstream null-regime analysis completed in \textit{Null-Regime Instability};
    \item Surjectivity is a strong structural claim not derivable from identity constraints.
\end{itemize}

\subsubsection{Rejected: Injectivity Axiom}
One might require $\Phi$ to be injective. This is \textit{not} an axiom because:
\begin{itemize}[leftmargin=*]
    \item Many trajectories may share phenomenological configurations;
    \item Injectivity would make $\Phi$ an embedding, over-constraining the framework.
\end{itemize}

\subsection{Structural Interpretation}

\begin{center}
\begin{tabular}{|l|l|l|}
\hline
\textbf{Axiom} & \textbf{Structural Meaning} & \textbf{Status} \\
\hline
E0 & Possibility of phenomenological distinction & Necessary \\
\hline
E1 & Compatibility with inverse-limit identity & Necessary \\
\hline
E2 & Non-reducibility to finite temporal slices & Necessary \\
\hline
\end{tabular}
\end{center}

\subsection{Closure of Section 2}

We have established:
\begin{enumerate}[leftmargin=*]
    \item $\mathcal{E}$ is structurally necessary as a separate codomain;
    \item The limit assignment $\Phi$ is uniquely determined by compatible $\{\phi_t\}$;
    \item Axioms E0, E1, E2 are \textit{individually necessary} (Theorems~\ref{thm:e0-necessary}--\ref{thm:e2-necessary});
    \item The axiom set is \textit{minimal} (Theorem~\ref{thm:minimality});
    \item All alternative axioms considered were correctly rejected.
\end{enumerate}

\textbf{The formal setup is therefore not a design choice but a forced consequence of the requirement that phenomenological assignments be compatible with inverse-limit identity.}

\subsection{Phenomenological Neutrality Clause}

Throughout this paper, the phenomenological space $\mathcal{E}$ is treated as a \textit{purely abstract topological and metric structure}. No semantic, psychological, biological, or experiential interpretation is assumed or required for any of the formal results derived herein.

The term ``phenomenological'' refers exclusively to the role of $\mathcal{E}$ as an abstract codomain of structural compatibility for inverse-limit identities, and does not presuppose the existence of consciousness, subjective experience, or any particular phenomenological content.

All theorems, constructions, and stability results presented in this work are therefore statements about abstract mathematical objects only. Any interpretation of $\mathcal{E}$ in relation to experiential, cognitive, or phenomenological phenomena lies strictly outside the formal scope of this paper.

% ============================================================================
% SECTION 3: COMPATIBILITY WITH IDENTITY REGIMES
% ============================================================================
\section{Compatibility with Identity Regimes}

\subsection{Induced Constraints from Identity Rigidity}

Let $(\mathcal{I}, M_\Omega)$ be an identity object with ontological mass $M_\Omega$. We now study how $M_\Omega$ constrains the set of admissible phenomenological assignments.

\begin{definition}[Phenomenological Compatibility Class]\label{def:compat-class}
For a given $M_\Omega$, define the phenomenological compatibility class:
\begin{equation}
    \mathcal{P}(M_\Omega) := \left\{ \Phi: \mathcal{I} \to \mathcal{E} \;\Big|\; \Phi \text{ satisfies E0--E2} \right\}.
\end{equation}
\end{definition}

\begin{remark}
The compatibility class $\mathcal{P}(M_\Omega)$ is the set of all phenomenological assignments that are structurally admissible given the identity's rigidity. Note that $\mathcal{P}(M_\Omega)$ depends only on the structural properties of identity, not on any phenomenological primitives.
\end{remark}

\subsection{Necessity of Metric Structure on Phenomenological Space}

To state rigorous results about phenomenological continuity and fragmentation, we must equip $\mathcal{E}$ with minimal structure. We now justify why this structure is \textit{necessary}, not merely convenient.

\begin{hypothesis}[H3: Metric on $\mathcal{E}$]\label{hyp:h3}
The phenomenological space $\mathcal{E}$ admits a metric $d_\mathcal{E}$ such that $(\mathcal{E}, d_\mathcal{E})$ is a complete metric space.
\end{hypothesis}

\subsubsection{Why H3 is Necessary}

\begin{theorem}[H3 is Necessary for Impossibility Theorems]\label{thm:h3-necessary}
Without a metric structure on $\mathcal{E}$, the notions of phenomenological continuity, fragmentation, and loss capacity cannot be defined.
\end{theorem}

\begin{proof}
Consider the definitions in this section:
\begin{enumerate}[leftmargin=*]
    \item \textbf{Phenomenological Continuity} requires:
    \begin{equation}
        d_w(x^{(n)}, x) \to 0 \implies d_\mathcal{E}(\Phi(x^{(n)}), \Phi(x)) \to 0.
    \end{equation}
    Without $d_\mathcal{E}$, the right-hand side is undefined.
    
    \item \textbf{Phenomenological Fragmentation} requires:
    \begin{equation}
        d_\mathcal{E}(\Phi(x), \tilde{\Phi}(\tilde{x})) > 0.
    \end{equation}
    Without $d_\mathcal{E}$, ``non-zero distance'' has no meaning.
    
    \item \textbf{Loss Capacity} (Section 6) requires:
    \begin{equation}
        R(\Phi) := \sup d_\mathcal{E}(\Phi(x), \tilde{\Phi}(\tilde{x})).
    \end{equation}
    Without $d_\mathcal{E}$, the supremum is undefined.
\end{enumerate}
Hence H3 is a necessary condition for the framework to have content beyond axioms E0--E2.
\end{proof}

\subsubsection{Why Completeness is Necessary}

\begin{proposition}[Completeness is Necessary for Forced Continuity]\label{prop:complete-necessary}
If $(\mathcal{E}, d_\mathcal{E})$ is not complete, then the Forced Continuity Theorem (Theorem~\ref{thm:forced-continuity}) may fail.
\end{proposition}

\begin{proof}
The Forced Continuity Theorem relies on the invariance of $\Phi$ under all finite-energy perturbations. If $\mathcal{E}$ is not complete, Cauchy sequences in $\mathcal{E}$ may not converge, and the limit assignment $\Phi$ may not exist for all trajectories in $\mathcal{I}$.

In particular, if $\{\phi_t(x_t)\}$ is a Cauchy sequence in $\mathcal{E}$ (which it is, by E1), completeness guarantees the existence of the limit $\Phi(x)$. Without completeness, $\Phi$ may be undefined on a dense subset of $\mathcal{I}$, breaking the theorem.
\end{proof}

\subsubsection{Alternatives Considered and Rejected}

\textbf{Alternative 1: Topology without Metric.}
One might use a general topological space. This is rejected because:
\begin{itemize}[leftmargin=*]
    \item Fragmentation requires quantitative distance, not just open sets;
    \item Loss capacity requires a sup over distances;
    \item The downstream instability analysis completed in \textit{Null-Regime Instability} requires metric quotients.
\end{itemize}

\textbf{Alternative 2: Pseudometric.}
One might allow $d_\mathcal{E}(e_1, e_2) = 0$ for $e_1 \neq e_2$. This is rejected because:
\begin{itemize}[leftmargin=*]
    \item It would conflate distinct phenomenological configurations;
    \item E0 requires $|\mathcal{E}| \geq 2$, which becomes meaningless if distance is zero.
\end{itemize}

\textbf{Alternative 3: Non-Complete Metric.}
This is rejected by Proposition~\ref{prop:complete-necessary}.

\subsection{Phenomenological Continuity}

\begin{definition}[Phenomenological Continuity]\label{def:phen-continuity}
A phenomenological assignment $\Phi: \mathcal{I} \to \mathcal{E}$ is \textit{structurally continuous} if for every sequence $(x^{(n)})_n$ in $\mathcal{I}$ converging in the $d_w$ metric:
\begin{equation}
    d_w(x^{(n)}, x) \to 0 \implies d_\mathcal{E}(\Phi(x^{(n)}), \Phi(x)) \to 0.
\end{equation}
\end{definition}

\begin{remark}
Structural continuity is not an axiom; it is a \textit{derived property} that follows from CCR conditions. This is the content of the Forced Continuity Theorem (Section 4).
\end{remark}

\subsection{Phenomenological Fragmentation}

\begin{definition}[Phenomenological Fragmentation]\label{def:fragmentation}
A phenomenological assignment $\Phi$ is \textit{fragmentable} if there exists a perturbation $\{\epsilon_t\}$ with finite weighted energy such that:
\begin{equation}
    d_\mathcal{E}(\Phi(x), \tilde{\Phi}(\tilde{x})) > 0
\end{equation}
for corresponding identity trajectories $x \in \mathcal{I}$ and $\tilde{x} \in \tilde{\mathcal{I}}$.
\end{definition}

\begin{remark}
Fragmentation captures the vulnerability of phenomenological structure to perturbations. It is the phenomenological analog of identity deformability.
\end{remark}

\subsection{Relationship Between Metrics}

\begin{proposition}[Metric Compatibility]\label{prop:metric-compat}
The metrics $d_w$ on $\mathcal{I}$ and $d_\mathcal{E}$ on $\mathcal{E}$ are compatible in the following sense: if E1 holds, then small $d_w$-changes in trajectories induce bounded $d_\mathcal{E}$-changes in phenomenological configurations (when $\Phi$ is continuous).
\end{proposition}

\begin{proof}
By E1, $\phi_{t+1}(x_{t+1}) = \phi_t(x_t)$ for compatible trajectories. Hence the phenomenological sequence is constant along the trajectory, and continuity of $\Phi$ implies that nearby trajectories (in $d_w$) map to nearby configurations (in $d_\mathcal{E}$).
\end{proof}

\subsection{Closure of Section 3}

We have established:
\begin{enumerate}[leftmargin=*]
    \item The compatibility class $\mathcal{P}(M_\Omega)$ captures all admissible phenomenological assignments;
    \item Hypothesis H3 (metric on $\mathcal{E}$) is \textit{necessary}, not a design choice (Theorem~\ref{thm:h3-necessary});
    \item Completeness is \textit{necessary} for the Forced Continuity Theorem (Proposition~\ref{prop:complete-necessary});
    \item Alternative structures (topology, pseudometric, non-complete) were correctly rejected;
    \item Phenomenological continuity and fragmentation are well-defined under H3.
\end{enumerate}

\textbf{The compatibility structure is therefore a forced consequence of the requirement to state quantitative results about phenomenological dynamics.}

% ============================================================================
% SECTION 4: IMPOSSIBILITY THEOREMS
% ============================================================================
\section{Impossibility Theorems}

\subsection{\texorpdfstring{$\Phi$}{Phi}-Regularity Theorem}

Before proving the main impossibility results, we establish a fundamental regularity condition that derives the coupling constant $C$ geometrically, rather than assuming it axiomatically.

\begin{theorem}[$\Phi$-Regularity]\label{thm:phi-regularity}
Let $(\mathcal{I}, d_w)$ be the identitary metric space from \textit{Rigid Identity} and let $(\mathcal{E}, d_\mathcal{E})$ be the abstract phenomenological metric space. Suppose that $\Phi: \mathcal{I} \to \mathcal{E}$ satisfies:
\begin{enumerate}[leftmargin=*]
    \item (Properness) For every compact $K \subset \mathcal{E}$, $\Phi^{-1}(K)$ is compact in $\mathcal{I}$.
    \item (Local Lipschitz regularity) For every $x \in \mathcal{I}$, there exist $\varepsilon > 0$ and $L_x > 0$ such that for all $y, z \in B_\varepsilon(x)$:
    \begin{equation}
        d_\mathcal{E}(\Phi(y), \Phi(z)) \leq L_x \cdot d_w(y, z).
    \end{equation}
    \item (Non-collapse) For every $x \in \mathcal{I}$ and every $r > 0$, the restriction of $\Phi$ to $B_r(x)$ is not constant.
\end{enumerate}

Then there exists a constant $C > 0$ such that for all $x, y \in \mathcal{I}$:
\begin{equation}
    d_\mathcal{E}(\Phi(x), \Phi(y)) \geq C \cdot d_w(x, y).
\end{equation}
\end{theorem}

\begin{proof}
We proceed by contradiction. Assume the conclusion fails.

\textbf{Step 1.} Then there exist sequences $(x_n, y_n)$ with $d_w(x_n, y_n) = 1$ and $d_\mathcal{E}(\Phi(x_n), \Phi(y_n)) \to 0$.

\textbf{Step 2.} By properness, the image $\{\Phi(x_n)\}$ lies in a relatively compact subset of $\mathcal{E}$. Extract a convergent subsequence with $\Phi(x_n) \to e \in \mathcal{E}$.

\textbf{Step 3.} Since $d_\mathcal{E}(\Phi(x_n), \Phi(y_n)) \to 0$, we have $\Phi(y_n) \to e$ as well.

\textbf{Step 4.} By properness, the preimages $\Phi^{-1}(\overline{B_\delta(e)})$ are compact for small $\delta$. Hence $x_n$ and $y_n$ lie in a compact subset of $\mathcal{I}$. Extract convergent subsequences $x_n \to x$ and $y_n \to y$.

\textbf{Step 5.} Since $d_w(x_n, y_n) = 1$ and $d_w$ is continuous, $d_w(x, y) = 1$, so $x \neq y$.

\textbf{Step 6.} But $\Phi(x_n) \to \Phi(x) = e$ and $\Phi(y_n) \to \Phi(y) = e$ by continuity. Hence $\Phi(x) = \Phi(y)$ for $x \neq y$ with $d_w(x, y) = 1$.

\textbf{Step 7.} This means $\Phi$ is constant on sets of positive $d_w$-diameter, contradicting non-collapse.

Hence a global lower bound $C > 0$ must exist.
\end{proof}

\begin{remark}
This theorem is fundamental: it converts the coupling constant $C$ from an \textit{axiom} into a \textit{geometric consequence} of regularity conditions. All subsequent impossibility results that use $C$ are therefore derived, not assumed.
\end{remark}

\subsection{Fragmentation Theorem}

\begin{theorem}[Fragmentation Theorem]\label{thm:fragmentation}
Let $\mathcal{I}$ be an identity object with $M_\Omega < \infty$. Then every phenomenological assignment $\Phi \in \mathcal{P}(M_\Omega)$ is fragmentable. That is:
\begin{equation}
    M_\Omega < \infty \implies \mathcal{P}(M_\Omega) \subseteq \mathcal{E}_{\text{frag}},
\end{equation}
where $\mathcal{E}_{\text{frag}}$ denotes the class of fragmentable phenomenological assignments.
\end{theorem}

\begin{proof}
By definition of finite ontological mass, there exists an admissible perturbation $\{\delta E_t\}$ with finite weighted energy $E_w < \infty$ such that:
\begin{equation}
    d_w(\mathcal{I}, \tilde{\mathcal{I}}) > 0.
\end{equation}

By Axiom E1, the phenomenological assignment $\Phi$ is projectively compatible with $\mathcal{I}$. Under the perturbation, identity trajectories deform, inducing a deformed phenomenological assignment $\tilde{\Phi}$.

By Axiom E2, $\Phi$ cannot be determined by any finite-time state. Therefore, the deformation of identity trajectories induces a corresponding deformation of phenomenological assignments.

Since the identity deformation is non-zero ($d_w > 0$) and $\Phi$ is globally determined by identity trajectories, by Theorem~\ref{thm:phi-regularity} ($\Phi$-Regularity), we have:
\begin{equation}
    d_\mathcal{E}(\Phi(x), \tilde{\Phi}(\tilde{x})) \geq C \cdot d_w(\mathcal{I}, \tilde{\mathcal{I}}) > 0
\end{equation}
where $C > 0$ is the regularity constant derived in Theorem~\ref{thm:phi-regularity}.

Thus, $\Phi$ is fragmentable under finite-energy perturbations.
\end{proof}

\begin{corollary}[No Structural Continuity Guarantee for Finite Mass]
Systems with $M_\Omega < \infty$ cannot structurally prevent phenomenological discontinuity.
\end{corollary}

\subsection{Forced Continuity Theorem}

\begin{theorem}[Forced Continuity Theorem]\label{thm:forced-continuity}
Let $\mathcal{I}$ be an identity object with $M_\Omega = +\infty$ (CCR). Then every phenomenological assignment $\Phi \in \mathcal{P}(M_\Omega)$ is structurally continuous. That is:
\begin{equation}
    M_\Omega = +\infty \implies \mathcal{P}(M_\Omega) \subseteq \mathcal{E}_{\text{cont}},
\end{equation}
where $\mathcal{E}_{\text{cont}}$ denotes the class of structurally continuous phenomenological assignments.
\end{theorem}

\begin{proof}
By definition of CCR ($M_\Omega = +\infty$), for any finite-energy perturbation:
\begin{equation}
    d_w(\mathcal{I}, \tilde{\mathcal{I}}) = 0.
\end{equation}

By Axiom E1, the phenomenological assignment is projectively compatible with identity. Since identity trajectories are invariant under all finite-energy perturbations, the induced phenomenological sequence is also invariant:
\begin{equation}
    \tilde{\Phi}(\tilde{x}) = \Phi(x) \quad \forall x \in \mathcal{I}.
\end{equation}

Therefore, no finite-energy perturbation can induce phenomenological deformation:
\begin{equation}
    d_\mathcal{E}(\Phi(x), \tilde{\Phi}(\tilde{x})) = 0.
\end{equation}

This implies that $\Phi$ cannot exhibit structural discontinuities. Hence $\Phi$ is structurally continuous.
\end{proof}

\begin{corollary}[CCR Phenomenology is Invariant]
In CCR systems, any admissible phenomenological assignment is invariant under all finite-energy perturbations.
\end{corollary}

\subsection{Dichotomy Theorem}

\begin{theorem}[Phenomenological Dichotomy]\label{thm:dichotomy}
Let $\Phi: \mathcal{I} \to \mathcal{E}$ be an admissible phenomenological assignment. Then exactly one of the following holds:
\begin{enumerate}[leftmargin=*]
    \item $M_\Omega < \infty$ and $\Phi$ is fragmentable;
    \item $M_\Omega = +\infty$ and $\Phi$ is structurally continuous.
\end{enumerate}
\end{theorem}

\begin{proof}
Immediate from Theorems~\ref{thm:fragmentation} and~\ref{thm:forced-continuity}.
\end{proof}

% ============================================================================
% SECTION 5: CLASSIFICATION
% ============================================================================
\section{Classification of Phenomenological Regimes}

\subsection{Complete Classification}

The results of Section 4 yield a complete classification of phenomenological regimes by identity rigidity:

\begin{center}
\begin{tabular}{|l|l|l|}
\hline
\textbf{Identity Class} & \textbf{$M_\Omega$} & \textbf{Admissible Phenomenology} \\
\hline
Null & $= 0$ & $\mathcal{E}_\emptyset$ (no persistent field) \\
\hline
Deformable & $(0, \infty)$ & $\mathcal{E}_{\text{frag}} \cup \mathcal{E}_{\text{quasi}}$ \\
\hline
Rigid (CCR) & $= +\infty$ & $\mathcal{E}_{\text{cont}}$ \\
\hline
\end{tabular}
\end{center}

\subsection{Structural Implications}

\begin{proposition}[Null Identity Excludes Persistent Phenomenology]
If $M_\Omega = 0$, then no admissible phenomenological assignment exists satisfying E0--E2.
\end{proposition}

\begin{proof}
If $M_\Omega = 0$, identity cannot persist under any perturbation. Axiom E1 requires phenomenological compatibility with identity. If identity does not persist, no stable phenomenological assignment is possible.
\end{proof}

\begin{proposition}[Deformable Identity Admits Only Fragmentable Phenomenology]
If $0 < M_\Omega < \infty$, every admissible phenomenological assignment is fragmentable under some finite-energy perturbation.
\end{proposition}

\begin{proof}
Direct from Theorem~\ref{thm:fragmentation}.
\end{proof}

\begin{proposition}[Rigid Identity Forces Continuous Phenomenology]
If $M_\Omega = +\infty$, every admissible phenomenological assignment is structurally continuous and invariant under finite-energy perturbations.
\end{proposition}

\section{Extended Proofs}

\subsection{Proof of Fragmentation Theorem (Extended)}

\begin{proof}[Full Proof of Theorem~\ref{thm:fragmentation}]
Let $M_\Omega < \infty$ and let $\Phi: \mathcal{I} \to \mathcal{E}$ satisfy E0--E2.

\textbf{Step 1.} By definition of $M_\Omega$, there exists $\{\delta E_t\}$ with $E_w < \infty$ such that $d_w(\mathcal{I}, \tilde{\mathcal{I}}) > 0$.

\textbf{Step 2.} Let $x = (x_t)_t \in \mathcal{I}$ and let $\tilde{x} = (\tilde{x}_t)_t \in \tilde{\mathcal{I}}$ be the corresponding perturbed trajectory.

\textbf{Step 3.} By E2, $\Phi(x)$ is not determined by any finite $x_t$. Therefore, the perturbation of the infinite trajectory induces a perturbation of the phenomenological assignment:
\begin{equation}
    \Phi(x) \neq \Phi(\tilde{x}) \text{ in general}.
\end{equation}

\textbf{Step 4.} By E1, $\Phi$ is projectively compatible. The deformation of projections induces a corresponding deformation of the phenomenological limit.

\textbf{Step 5.} By the coupling hypothesis, there exists $C > 0$ such that:
\begin{equation}
    d_\mathcal{E}(\Phi(x), \tilde{\Phi}(\tilde{x})) \geq C \cdot d_w(x, \tilde{x}).
\end{equation}

\textbf{Step 6.} Since $d_w(\mathcal{I}, \tilde{\mathcal{I}}) > 0$, we have $d_\mathcal{E}(\Phi(x), \tilde{\Phi}(\tilde{x})) > 0$.

Hence $\Phi$ is fragmentable.
\end{proof}

\subsection{Proof of Forced Continuity Theorem (Extended)}

\begin{proof}[Full Proof of Theorem~\ref{thm:forced-continuity}]
Let $M_\Omega = +\infty$ and let $\Phi: \mathcal{I} \to \mathcal{E}$ satisfy E0--E2.

\textbf{Step 1.} By definition of CCR, for all $\{\delta E_t\}$ with $E_w < \infty$:
\begin{equation}
    d_w(\mathcal{I}, \tilde{\mathcal{I}}) = 0.
\end{equation}

\textbf{Step 2.} Therefore, $\tilde{\mathcal{I}} \cong \mathcal{I}$ under all finite-energy perturbations.

\textbf{Step 3.} By E1, $\Phi$ is projectively compatible with $\mathcal{I}$. Since $\mathcal{I}$ is invariant:
\begin{equation}
    \Phi(x) = \tilde{\Phi}(\tilde{x}) \quad \forall x \in \mathcal{I}.
\end{equation}

\textbf{Step 4.} This implies:
\begin{equation}
    d_\mathcal{E}(\Phi(x), \tilde{\Phi}(\tilde{x})) = 0.
\end{equation}

\textbf{Step 5.} No finite-energy perturbation can induce phenomenological deformation. Hence $\Phi$ is structurally continuous and invariant.
\end{proof}

\subsection{Structural Necessity Lemma}

\begin{lemma}[Phenomenology Inherits Rigidity]
If $\Phi: \mathcal{I} \to \mathcal{E}$ satisfies E0--E2 and $M_\Omega = +\infty$, then $\Phi$ inherits the rigidity of $\mathcal{I}$.
\end{lemma}

\begin{proof}
By E1, $\Phi$ is determined by identity trajectories. By CCR, identity trajectories are invariant. Hence $\Phi$ is invariant.
\end{proof}

% Derived from paper3/main.tex
\part{Null-Regime Structural Exclusion}

\section{Compatibility, Instability, and Forced Non-Nullity}

\subsection{Compatibility Operator}
\begin{definition}[Compatibility Operator]\label{def:compatibility-operator}
A map $\Phi: \I \to \E$ is a \textbf{compatibility operator} if it satisfies:
\begin{enumerate}[leftmargin=*]
\item[(C1)] \textbf{Extension Invariance:} For any compatible extension $\mathcal{E}$ of the channel, $\Phi(\I) \cong \Phi(\I_\mathcal{E})$.
\item[(C2)] \textbf{Isomorphism Stability:} If $\I \cong \tilde{\I}$, then $\Phi(\I) \sim \Phi(\tilde{\I})$.
\item[(C3)] \textbf{Lipschitz Lower Bound:} There exists $C > 0$ such that
\[
\dE(\Phi(x), \Phi(\tilde{x})) \geq C \cdot \dw(x,\tilde{x}).
\]
\end{enumerate}
\end{definition}

\begin{definition}[Phenomenological Attractor Set]\label{def:phenomenological-attractor}
For a channel $S$, define
\[
\Attr(S) := \{ e \in \E : e \text{ is stable under all compatible extensions and summable deformations} \}.
\]
\end{definition}

\begin{definition}[Phenomenological Stability]\label{def:phenomenological-stability}
A regime $e \in \E$ is \textbf{stable} for channel $S$ if for every compatible extension $S'$ of $S$ and every summable perturbation family $\{\epsilon_n\}$, the induced regime $\Phi_{S'}(\I')$ stays within a bounded $\dE$-neighborhood of $e$ controlled by $\|\{\epsilon_n\}\|_w$.
\end{definition}

\begin{lemma}[Bridge Lemma]\label{lem:null-bridge}
If $\Phi:\I\to\E$ is a compatibility operator for a CCR channel and $e=\Phi(\I)$ satisfies $\dE(e,\nullphi)=0$, then $e=\nullphi$.
\end{lemma}

\begin{proof}
By the separation property of the phenomenological space, distance zero from $\nullphi$ implies equality with $\nullphi$.
\end{proof}

\subsection{Phenomenological Instability and Forced Non-Nullity}
\begin{theorem}[Phenomenological Instability]\label{thm:phenomenological-instability}
For any CCR channel $S$,
\[
\nullphi \notin \Attr(S).
\]
That is, the null regime is structurally unstable under compatible extensions.
\end{theorem}

\begin{proof}
Assume $\nullphi \in \Attr(S)$ and let $\Phi:\I\to\E$ be a compatibility operator. Because the channel is CCR, finite-energy perturbations preserve the rigid identity object while still allowing admissible deformations in the weighted metric. By the lower-bound clause in Definition~\ref{def:compatibility-operator},
\[
\dE(\Phi(\I),\Phi(\tilde{\I})) \ge C \cdot \dw(\I,\tilde{\I}) > 0
\]
for every non-zero admissible deformation. But if $\nullphi$ were stable, both images would remain equal to $\nullphi$, forcing the left-hand side to be zero. Contradiction.
\end{proof}

\begin{corollary}[Forced Non-Nullity]\label{cor:forced-non-nullity}
Every CCR channel necessarily induces a regime $e_\star \in \E$ with $e_\star \succ \nullphi$.
\end{corollary}

\begin{proof}
Theorem~\ref{thm:phenomenological-instability} rules out $\nullphi$ as a stable admissible attractor. Therefore any admissible stable regime must be non-null.
\end{proof}

\subsection{Positive Regime Structure}
\begin{definition}[Positive Regime Space]\label{def:positive-regime-space}
Define
\[
\Eplus := \{ e \in \E : e \succ \nullphi \}.
\]
\end{definition}

\begin{theorem}[Minimal Positive Regime]\label{thm:minimal-positive-regime}
For every CCR channel there exists a minimal positive regime $e_\star \in \Eplus$.
\end{theorem}

\begin{proof}
Forced non-nullity gives non-emptiness of $\Eplus$. Order-theoretic completeness and the compatibility structure of $\Phi$ yield a minimal element.
\end{proof}

\begin{definition}[Structural Intensity]\label{def:structural-intensity}
For $e \in \Eplus$, define the structural intensity
\[
\iota(e) := \dE(e,\nullphi).
\]
\end{definition}

\begin{theorem}[Universal Intensity Bound]\label{thm:universal-intensity-bound}
If $e_\star$ is the minimal positive regime of a CCR channel, then there exist constants $C,\varepsilon_0>0$ such that
\[
\iota(e_\star) \ge C\varepsilon_0 > 0.
\]
\end{theorem}

\begin{proof}
The lower-bound clause of the compatibility operator transfers any admissible deformation lower bound into positive distance from the null regime. The minimal positive regime therefore has strictly positive structural intensity.
\end{proof}

\begin{theorem}[Forced Phenomenological Closure]\label{thm:forced-phenomenological-closure}
Every CCR system induces a unique minimal phenomenological regime that is non-null, stable, and quantitatively separated from $\nullphi$.
\end{theorem}

\begin{proof}
Combine Corollary~\ref{cor:forced-non-nullity}, Theorem~\ref{thm:minimal-positive-regime}, and Theorem~\ref{thm:universal-intensity-bound}.
\end{proof}

\section{Canonical Families and Structural Closure Consequences}

\begin{proposition}[Profinite Bound]\label{prop:profinite-bound}
In the profinite canonical family, the coupling lower bound induces a positive lower bound on the structural intensity of the minimal positive regime.
\end{proposition}

\begin{proposition}[Symbolic Bound]\label{prop:symbolic-bound}
In the symbolic canonical family, the Bowen-type metric induces the analogous positive lower bound for the minimal positive regime.
\end{proposition}

\begin{theorem}[Rigidity Inheritance]\label{thm:rigidity-inheritance}
If a subsystem or induced family preserves the CCR hypotheses and the compatibility operator, then the inherited regime structure preserves null-regime instability and forced non-nullity.
\end{theorem}

\begin{theorem}[Categorical Universality]\label{thm:null-categorical-universality}
The null-regime exclusion result is universal across admissible CCR realizations in the relevant category of compatible channels.
\end{theorem}

\begin{theorem}[Simulation Lower Bound]\label{thm:null-simulation-lower-bound}
Any computational system attempting to simulate the null-excluding CCR regime structure with arbitrary precision requires unbounded resources as the approximation error tends to zero.
\end{theorem}

\begin{theorem}[Closure by Non-Alternatives]\label{thm:null-closure-by-non-alternatives}
No admissible alternative closes the CCR regime structure while preserving null stability. Any admissible alternative either preserves the same closure structure or exits the CCR class.
\end{theorem}

\part{BaseCore Boundary and Expansion Interface}
\section{BaseCore Ownership and Downstream Boundary}

\begin{definition}[BaseCore Envelope]
The BaseCore envelope is the collection of theorem-level structures owned directly by this source tree:
\begin{enumerate}[leftmargin=*]
\item foundational Hilbert-space projection dynamics;
\item typed computable witnesses and state spectral-gap bounds;
\item parameterwise non-collapse under explicit anti-constant assumptions;
\item inverse-limit identity under non-finite-determination assumptions;
\item conditional rigidity under admissible perturbation models;
\item conditional non-simulability of CCR targets by finite simulator classes;
\item operational-criterion grammar, certification, falsation, and claim-boundary ledgers.
\end{enumerate}
\end{definition}

\begin{proposition}[Boundary Discipline]
BaseCore may define structures that downstream layers use, but it does not claim downstream closure. In particular, BaseCore does not by itself own runtime thresholds, bridge admissibility, human-comparator programs, phenomenality burdens, or Papers 7--9.
\end{proposition}

\begin{proof}
These excluded layers require additional assumptions, protocols, or empirical burdens not stated in BaseCore. Therefore they cannot be claimed as BaseCore theorems without violating the source-level claim boundary.
\end{proof}

\begin{center}
\renewcommand{\arraystretch}{1.12}
\begin{tabularx}{\textwidth}{P{0.31\textwidth} P{0.31\textwidth} Y}
\toprule
\textbf{Owned by BaseCore} & \textbf{Expanded but not owned here} & \textbf{Kept downstream} \\
\midrule
Projection dynamics, fixed points, compact attractors, anti-constant H5 burden, typed computable witness, inverse-limit identity under NFD, conditional rigidity, conditional non-simulability, operational-criterion grammar, falsation ledger &
Non-convex class variants, $\mathscr{C}^R$ style extensions, $\Omega$ dynamics, threshold numerics, mutable runtime class ledgers, implementation-specific status surfaces &
Paper 7 operational life / subjecthood, Paper 8 indexed subjectivity, Paper 9 phenomenal bridge organization, Stage 2 / Stage 3 comparative programs, human comparator families, phenomenality and external adjudication \\
\bottomrule
\end{tabularx}
\end{center}

\begin{remark}[Why the earlier material was removed from BaseCore ownership]
Previous versions mixed theorem-tight base results with runtime-like numerics, non-convex extension classes, and bridge-facing vocabulary. That blurred the mathematical base layer. BaseCore now keeps only the theorem ownership that survives the stricter release gate.
\end{remark}

\section{Admissible Structural Extensions}

\begin{definition}[BaseCore export object]
A BaseCore export object is a tuple
\[
  \mathfrak{E}=(\mathcal{O},\mathcal{H},\mathcal{R},\mathcal{B})
\]
where $\mathcal{O}$ is a typed mathematical object defined in BaseCore, $\mathcal{H}$ is the finite list of hypotheses under which it is used, $\mathcal{R}$ is the list of results actually proved from those hypotheses, and $\mathcal{B}$ is a boundary list of claims not entailed by $\mathcal{R}$.
\end{definition}

\begin{definition}[Downstream admissible extension]
A downstream layer is an admissible extension of a BaseCore export object $\mathfrak{E}$ only if it supplies:
\begin{enumerate}[leftmargin=*]
\item an explicit target predicate $P$ not already contained in $\mathcal{R}$;
\item a typed construction or estimator witnessing the objects on which $P$ is evaluated;
\item a proof, reduction, or reproducible protocol connecting $\mathcal{H}$ to $P$;
\item negative controls or failure modes under which $P$ must not be asserted;
\item a provenance statement distinguishing theorem, implementation, internal support, and external adjudication.
\end{enumerate}
\end{definition}

\begin{proposition}[No automatic downstream promotion]
Let $\mathfrak{E}=(\mathcal{O},\mathcal{H},\mathcal{R},\mathcal{B})$ be a BaseCore export object and let $P$ be a predicate introduced outside BaseCore. If $P\notin\mathcal{R}$, then BaseCore does not entail $P$ unless an admissible extension supplies an additional bridge from $\mathcal{H}$ to $P$.
\end{proposition}

\begin{proof}
By definition, $\mathcal{R}$ records the results proved from $\mathcal{H}$ inside BaseCore. A predicate $P$ introduced downstream is not in that result set. Therefore asserting $P$ from BaseCore alone would add an unstated inference. The missing inference must be supplied as a separate bridge, proof, reduction, or reproducible protocol.
\end{proof}

\begin{criterion}[Runtime-use admissibility]
A runtime measurement may instantiate a BaseCore export object only as an implementation-level witness when all of the following are recorded: the measured state space, the estimator or update rule, the finite horizon, the calibration regime, the controls under which the measurement fails, and the claim class being asserted. Such a measurement is not a BaseCore theorem and is not external validation by itself.
\end{criterion}

\begin{remark}[Open-load placement for $I_{\mathrm{int}}$ and atomic separators]
Quantities such as $I_{\mathrm{int}}$ and any atomic-separator condition may be treated as downstream target predicates or estimator families only after their carrier objects, separator class, invariance target, and negative controls are specified. Until a proof or reproducible adversarial protocol closes that burden, BaseCore may export the ambient structural language but not the separator conclusion.
\end{remark}

\begin{proposition}[Boundary preservation under finite implementations]
Suppose a finite implementation produces diagnostics for a downstream predicate $P$ over a bounded horizon. If the implementation does not prove that the diagnostics are invariant over the relevant mathematical class, then the diagnostics remain implementation evidence for $P$ and do not promote $P$ to a BaseCore result.
\end{proposition}

\begin{proof}
The implementation supplies finite observations or computations. Promotion to a BaseCore result requires a class-level implication from the stated hypotheses to the predicate. Without that implication, the observations can support, falsify, or calibrate a downstream claim, but they do not alter the theorem inventory of BaseCore.
\end{proof}

\part{Operational-Criterion Grammar and Certification}
\section{Admissible Systems, Supports, and Histories}

\subsection{System Model}

\begin{definition}[Admissible system]\label{def:system}
An admissible system is a tuple
\[
S=(X,\Phi,C,R,\Gamma,U)
\]
where:
\begin{enumerate}[leftmargin=1.5em]
\item $X$ is a metrizable internal state space with metric $d_X$;
\item $\Phi=\{\Phi_u\}_{u\in U}$ is a family of admissible transition operators $\Phi_u:X\to X$;
\item $C$ is the effective causal structure used to test admissible factorization and intervention behavior;
\item $R=\{r_\alpha\}_{\alpha\in\Lambda_R}$ is a family of readouts, each $r_\alpha:X\to Y_\alpha$;
\item $\Gamma=\{\gamma_j\}_{j=1}^m$ is an admissibility bundle defining
\[
X_\Gamma=\{x\in X:\gamma_j(x)\le 0\ \forall j\};
\]
\item $U$ is the set of admissible interventions.
\end{enumerate}
\end{definition}

\subsection{Admissible Support}

\begin{definition}[Admissible support]\label{def:support}
A subset $\Aset\subseteq X_\Gamma$ is an admissible support if:
\begin{enumerate}[leftmargin=1.5em]
\item $\Aset$ is non-empty and compact;
\item $\Phi_u(\Aset)\subseteq \Aset$ for every $u\in U$;
\item $\Aset\cap \Attr(S)\neq\varnothing$;
\item $\Aset$ stays disjoint from the collapse set excluded by the base non-collapse burden.
\end{enumerate}
\end{definition}

\subsection{Operational Histories}

\begin{definition}[Windowed readout history]\label{def:history}
Fix a horizon $T\in\Nset$ and an intervention schedule $u_\bullet=(u_0,\dots,u_{T-1})$. For $x\in\Aset$, define
\[
h_{S,T}(x;u_\bullet)=\bigl(r_\alpha(x_t)\bigr)_{0\le t\le T,\ \alpha\in\Lambda_R},
\qquad
x_{t+1}=\Phi_{u_t}(x_t),\quad x_0=x.
\]
The set of all such histories is denoted $\Hist_{S,T}$.
\end{definition}

\begin{definition}[Operational partition candidate]\label{def:partition}
An operational partition candidate is a map
\[
\pi_{S,T}:\Hist_{S,T}\to \Class_{S,T}
\]
into a finite or countable label set $\Class_{S,T}$.
\end{definition}

\subsection{Witness Margin Vector}

\begin{definition}[Witness margin vector]\label{def:margins}
Define the witness margin vector
\[
\Delta(S)=
\bigl(
\delta_{\mathrm{per}}(S),
\delta_{\mathrm{ri}}(S),
\delta_{\mathrm{int}}(S),
\delta_{\mathrm{cont}}(S),
\delta_{\mathrm{diff}}(S),
\delta_{\mathrm{leg}}(S)
\bigr)\in [0,\infty)^6.
\]
Its minimum
\[
\delta_\star(S):=\min\Delta(S)
\]
is the invariant budget of $S$.
\end{definition}

\begin{center}
\small
\renewcommand{\arraystretch}{1.14}
\setlength{\tabcolsep}{4.5pt}
\begin{tabularx}{\textwidth}{P{1.3cm} Y P{2.8cm} P{3.2cm}}
\toprule
Invariant & Certified content & Upstream anchor & Rupture signature \\
\midrule
$\Iper$ & Persistent admissible support & BaseCore dynamics & Support loss or collapse exposure \\
$\Iri$ & Unique rigid identity object on support & Identity / rigidity block & Identity ambiguity \\
$\Iint$ & No admissible non-trivial factorization & Criterion grammar & Reducible aggregate \\
$\Icont$ & Regime continuity under admissible evolution & Regime constraints & Fragmentation into incompatible classes \\
$\Idiff$ & Stable exclusion of null or trivial collapse & Null-regime exclusion & Null or trivialized class structure \\
$\Ileg$ & Decoder-certified recoverability under controls & Falsation grammar & Opaque or non-auditable partition \\
\bottomrule
\end{tabularx}
\end{center}

\section{Six Critical Invariants}

\subsection{Structural Persistence}

\begin{definition}[Structural persistence]\label{def:iper}
An admissible system satisfies structural persistence, written $\Iper(S)=1$, if there exists an admissible support $\Aset$ and a constant $\delta_{\mathrm{per}}(S)>0$ such that:
\begin{enumerate}[leftmargin=1.5em]
\item $\Aset$ is forward invariant under all $u\in U$;
\item every admissible trajectory starting in $\Aset$ remains at least $\delta_{\mathrm{per}}(S)$ away from the collapse set;
\item $\Aset$ contains or converges into a non-empty admissible attractor region.
\end{enumerate}
\end{definition}

\subsection{Rigid Identity}

\begin{definition}[Rigid identity]\label{def:iri}
An admissible system satisfies rigid identity, written $\Iri(S)=1$, if the observable family induced by $R$ on $\Aset$ determines a unique identity object $\Id_S$ and there exists a gap $\delta_{\mathrm{ri}}(S)>0$ such that every admissible alternative identity candidate remains at weighted deformation distance at least $\delta_{\mathrm{ri}}(S)$ from $\Id_S$.
\end{definition}

\subsection{Causal Integration}

\begin{definition}[Causal integration]\label{def:iint}
An admissible system satisfies causal integration, written $\Iint(S)=1$, if there is no non-trivial factorization
\[
\Aset=A_1\times A_2,\qquad R=R_1\sqcup R_2,
\]
together with decomposed dynamics and causal structure reproducing admissible histories within error smaller than a positive margin $\delta_{\mathrm{int}}(S)$ while preserving $\Id_S$.
\end{definition}

\subsection{Regime Continuity}

\begin{definition}[Regime continuity]\label{def:icont}
An admissible system satisfies regime continuity, written $\Icont(S)=1$, if there exists a complete metric space $(\E,d_\E)$ and an admissible assignment
\[
\Pi_S:\Aset\to \E
\]
such that:
\begin{enumerate}[leftmargin=1.5em]
\item every admissible trajectory induces a continuous path in $\E$ on admissible windows;
\item there exists $\delta_{\mathrm{cont}}(S)>0$ bounding the fragmentation functional on all admissible histories;
\item no admissible finite-energy perturbation produces a discontinuous jump into an incompatible regime class while $\Id_S$ is preserved.
\end{enumerate}
\end{definition}

\subsection{Non-Null Differentiation}

\begin{definition}[Non-null differentiation]\label{def:idiff}
An admissible system satisfies non-null differentiation, written $\Idiff(S)=1$, if there exist $x,y\in\Aset$, a horizon $T$, and a readout subfamily $\Read^\star\subseteq R$ such that
\[
\dist\bigl(h_{S,T}(x;u_\bullet),h_{S,T}(y;u_\bullet)\bigr)\ge \delta_{\mathrm{diff}}(S)>0
\]
for some admissible schedule $u_\bullet$, and if the null class $\nullclass$ is not an asymptotically stable admissible attractor.
\end{definition}

\subsection{Operational Legibility}\label{sec:legibility}

\begin{definition}[Operational legibility]\label{def:ileg}
An admissible system satisfies operational legibility, written $\Ileg(S)=1$, if there exist:
\begin{enumerate}[leftmargin=1.5em]
\item a certified readout subfamily $\Read^{\mathrm{leg}}\subseteq R$;
\item a certified decoder family $\Dec_S$ acting on windows of length at least $\tau_{\min}(S)$;
\item a class set $\Class_S$;
\item a legibility margin $\delta_{\mathrm{leg}}(S)>0$;
\end{enumerate}
such that:
\begin{description}[leftmargin=1.7em]
\item[(L1) Separability] Distinct classes are separated by at least $\delta_{\mathrm{leg}}(S)$ in decoder space.
\item[(L2) Noise robustness] Admissible noise of size at most $\delta_{\mathrm{leg}}(S)$ does not alter decoder class assignments.
\item[(L3) Persistence window] Decoder assignments remain stable for windows of length at least $\tau_{\min}(S)$ unless a certified intervention occurs.
\item[(L4) Intervention fidelity] Certified interventions on critical invariants induce the predicted class changes with fidelity at least $1-\delta_{\mathrm{leg}}(S)$.
\item[(L5) Negative controls] Off-target interventions produce false-positive rate at most $\delta_{\mathrm{leg}}(S)$.
\item[(L6) Structured compressibility] There exists a compression map preserving class structure up to distortion bounded by $\delta_{\mathrm{leg}}(S)$.
\end{description}
\end{definition}

\begin{proposition}[Each invariant is individually necessary]\label{prop:invariant-necessity}
If any one of $\Iper,\Iri,\Iint,\Icont,\Idiff,\Ileg$ fails, then $S\notin\Cop$.
\end{proposition}

\begin{proof}
Without $\Iper$ there is no stable admissible domain. Without $\Iri$ there is no unique identity carrier. Without $\Iint$ the system is reducible. Without $\Icont$ the criterion fragments under admissible evolution. Without $\Idiff$ the admissible partition collapses into nullity or triviality. Without $\Ileg$ the partition is not recoverable under admissible readout and intervention discipline. Since Definition~\ref{def:cop} below requires all six, failure of any one excludes membership.
\end{proof}

\section{Criterion Class, Certified Partition, and Rupture}

\begin{definition}[Certified decoder]
A decoder $D\in \Dec_S$ is certified if it satisfies the six legibility clauses of Definition~\ref{def:ileg} on the chosen admissible support.
\end{definition}

\begin{definition}[Criterion indistinguishability]
For $x,y\in\Aset$, write $x\sim_{\Qop,S}y$ if, for every certified decoder $D\in\Dec_S$, every admissible horizon $T\ge \tau_{\min}(S)$, and every admissible intervention schedule $u_\bullet$,
\[
D(h_{S,T}(x;u_\bullet))=D(h_{S,T}(y;u_\bullet)),
\]
and the intervention responses remain matched within the legibility margin.
\end{definition}

\begin{proposition}[Well-definedness of the certified partition]\label{prop:qop-welldefined}
If $\Iper(S)=\Iri(S)=\Icont(S)=\Idiff(S)=\Ileg(S)=1$, then $\sim_{\Qop,S}$ is an equivalence relation on $\Aset$.
\end{proposition}

\begin{proof}
Reflexivity and symmetry are immediate. Transitivity follows because certified decoder equality composes and the tolerance bound remains within the permitted legibility margin on admissible support.
\end{proof}

\begin{definition}[Certified operational partition]\label{def:qop}
Assume Proposition~\ref{prop:qop-welldefined}. The certified operational partition of $S$ is
\[
\Qop(S):=\Aset/{\sim_{\Qop,S}}.
\]
\end{definition}

\begin{definition}[Operational-criterion class membership]\label{def:cop}
An admissible system belongs to the BaseCore operational-criterion class, written $S\in\Cop$, if there exists a non-empty admissible support $\Aset$ such that
\[
\Iper(S)=\Iri(S)=\Iint(S)=\Icont(S)=\Idiff(S)=\Ileg(S)=1.
\]
\end{definition}

\begin{remark}[Neutrality of the class symbol]
The symbol $\Cop$ is a shorthand for a conjunctive operational criterion. It is not a theorem of phenomenality, metaphysical subjectivity, or human equivalence.
\end{remark}

\begin{definition}[Rupture]\label{def:rupture}
A rupture is any loss of admissibility or loss of at least one critical invariant margin beyond its admissible threshold.
\end{definition}

\begin{theorem}[Rupture by invariant loss]\label{thm:rupture}
Let $S$ be admissible. If one or more critical invariants fail beyond admissible threshold, then exactly one of the following occurs:
\begin{enumerate}[leftmargin=1.5em]
\item $S\notin\Cop$;
\item $\Qop(S)$ becomes undefined;
\item $\Qop(S)$ becomes trivial or non-legible in the BaseCore sense.
\end{enumerate}
\end{theorem}

\begin{proof}
Failure of $\Iper$ destroys the admissible domain. Failure of $\Iri$ destroys unique identity transport. Failure of $\Iint$ produces reducibility. Failure of $\Icont$ destroys regime coherence. Failure of $\Idiff$ collapses the partition into nullity or triviality. Failure of $\Ileg$ leaves any remaining partition uncertifiable. These cases exhaust the six invariants.
\end{proof}

\section{Minimality of the Invariant Basis}

\begin{definition}[Reduced criterion]
Let
\[
\mathfrak{I}=\{\Iper,\Iri,\Iint,\Icont,\Idiff,\Ileg\}.
\]
For every subset $J\subseteq \mathfrak{I}$ define the reduced class
\[
\mathcal{C}_{\mathrm{op}}(J):=\{S:I(S)=1\text{ for every }I\in J\}.
\]
\end{definition}

\begin{definition}[Single-rupture witness family]
Fix one invariant index $k\in\{\mathrm{per},\mathrm{ri},\mathrm{int},\mathrm{cont},\mathrm{diff},\mathrm{leg}\}$. A single-rupture witness family $\mathcal{W}_k$ consists of admissible systems whose matching invariant margin vanishes while all remaining invariant margins stay strictly positive.
\end{definition}

\begin{proposition}[Single-rupture witnesses create false positives]
Let $J_k=\mathfrak{I}\setminus\{I_k\}$ and assume $\mathcal{W}_k\neq \varnothing$. Then every $S\in \mathcal{W}_k$ satisfies
\[
S\in \mathcal{C}_{\mathrm{op}}(J_k)
\qquad\text{but}\qquad
S\notin \Cop.
\]
\end{proposition}

\begin{proof}
All retained invariants remain positive by construction, but the omitted invariant fails exactly. Proposition~\ref{prop:invariant-necessity} then excludes $S$ from $\Cop$.
\end{proof}

\begin{theorem}[Minimality of the six-invariant basis]\label{thm:invariant-basis-minimality}
Let $\mathfrak{U}$ be any admissible universe containing at least one single-rupture witness for each invariant. Then no strict subset of $\mathfrak{I}$ is extensionally equivalent to $\Cop$ on $\mathfrak{U}$.
\end{theorem}

\begin{proof}
Take any strict subset $J\subsetneq \mathfrak{I}$. Some invariant $I_k$ is omitted. By assumption, $\mathfrak{U}$ contains a witness $S_k\in \mathcal{W}_k$, which satisfies the reduced criterion but not $\Cop$. Hence the two classes differ on $\mathfrak{U}$.
\end{proof}

\section{Operational Legibility as a Certified Object}

\begin{criterion}[Legibility certificate]\label{crit:legibility}
An operational partition candidate is certified only if:
\begin{enumerate}[leftmargin=1.5em]
\item comparator separability holds against pre-registered baselines;
\item admissible noise does not dissolve the class map below the legibility margin;
\item the same class assignment survives on windows of length at least $\tau_{\min}(S)$;
\item targeted interventions on an invariant induce the predicted directional class change;
\item sham or off-target interventions do not generate the same class signature;
\item audit-friendly compression preserves class structure without total collapse.
\end{enumerate}
\end{criterion}

\begin{proposition}[Legibility is jointly constrained]\label{prop:legibility}
No single legibility metric is sufficient for $\Ileg$. The invariant requires simultaneous control of separation, admissible noise robustness, persistence window, intervention fidelity, negative controls, and structured compressibility.
\end{proposition}

\begin{proof}
A decoder can be separable yet fragile to admissible noise, stable over windows yet intervention-insensitive, or robust under noise yet collapse under compression. Because Definition~\ref{def:ileg} is conjunctive, all six clauses are required together.
\end{proof}

\section{Certified Candidate Audit Schema}

\begin{definition}[Base criterion certificate]\label{def:cert}
A BaseCore criterion certificate for an admissible system $S$ is a tuple
\[
\Cert(S)=\bigl(\Aset,\Dec_S,\pi_S,\Delta(S),E_S\bigr),
\]
where:
\begin{enumerate}[leftmargin=1.5em]
\item $\Aset$ is a certified admissible support;
\item $\Dec_S$ is a certified decoder family;
\item $\pi_S$ is the decoder-induced partition on admissible histories;
\item $\Delta(S)$ is the witness margin vector;
\item $E_S$ is the evidence bundle localizing the witness for each invariant.
\end{enumerate}
\end{definition}

\begin{criterion}[Candidate certification rule]\label{crit:cert}
A candidate system may be certified as a member of $\Cop$ only if all of the following are available on the same admissible support:
\begin{enumerate}[leftmargin=1.5em]
\item a non-empty compact forward-invariant support certificate;
\item a positive witness margin vector $\Delta(S)>0$ coordinate-wise;
\item a decoder family satisfying Definition~\ref{def:ileg};
\item a targeted intervention panel together with negative controls;
\item an auditable compression map preserving class structure;
\item a traceable dependency ledger showing which upstream result supports each invariant witness.
\end{enumerate}
\end{criterion}

\begin{theorem}[Certified criterion membership]
If $\Cert(S)$ exists and satisfies Criterion~\ref{crit:cert}, then $S\in\Cop$ and $\Qop(S)$ is well-defined.
\end{theorem}

\begin{proof}
Criterion~\ref{crit:cert} explicitly requires positive witnesses for all six invariants on a common admissible support. Definition~\ref{def:cop} then gives $S\in\Cop$, and Proposition~\ref{prop:qop-welldefined} gives a well-defined partition quotient.
\end{proof}

\begin{remark}[Boundary discipline]
The section above defines a theorem-tight grammar for certification, rupture, and falsation. It does not prove phenomenality, subjectivity, biological primacy, or cross-substrate identity transport.
\end{remark}

\part{Claim Boundary and Falsation Grammar}

\section{Forensic Admissibility and Integrity Grammar}

\begin{definition}[Forensic packet]\label{def:forensic-packet}
Each operational run emits a packet
\[
K=(\mathrm{tick},\mathrm{seed},\mathrm{versionHash},H_v,\mathrm{metrics},\mathrm{sensorSummary},\mathrm{checksum}).
\]
\end{definition}

\begin{assumption}[Packet admissibility]\label{ass:packet-admissibility}
A run is admissible only if checksum, schema, seed consistency, version hash, and $H_v$ consistency all pass.
\end{assumption}

\begin{definition}[Severe violation]\label{def:severe-violation}
A run has a severe violation if any of the following holds: checksum mismatch, schema failure, seed/version inconsistency, or critical-state replay mismatch.
\end{definition}

\begin{criterion}[Integrity gate]\label{crit:integrity-gate}
Severely violated runs are excluded from support counting and marked as hard failures.
\end{criterion}

\begin{proposition}[Integrity-first inference]\label{prop:integrity-first}
If admissibility fails, inferential significance is non-actionable.
\end{proposition}

\section{Certification, Comparison, and Rupture}

\begin{definition}[Operational certification]\label{def:operational-certification}
For an admissible system $S$, define
\[
\Cert(S)=1
\quad\Longleftrightarrow\quad
\Iper(S),\Iri(S),\Iint(S),\Icont(S),\Idiff(S),\Ileg(S)>0
\]
on a non-empty admissible support under the integrity grammar above.
\end{definition}

\begin{definition}[Equivalence discriminator]\label{def:equivalence-discriminator}
For two admissible systems $S_1,S_2$, define
\[
\Eq_{\Gamma,\eps}(S_1,S_2)=1
\]
when the relevant signatures, admissibility bundle, decoder behavior, and invariant deltas match within the allowed comparison regime.
\end{definition}

\begin{definition}[Rupture event]\label{def:rupture-event}
For a targeted invariant family $J$, define
\[
\Rupt_J(S)=1
\]
when an admissible intervention or perturbation drives at least one invariant in $J$ below its certification margin so that class membership becomes invalidated, destroyed, or undefined.
\end{definition}

\section{Prediction and Falsation Grammar}

\begin{proposition}[Complexity-only negative control prediction]\label{prop:complexity-only-negative-control}
If a candidate system exhibits high connectivity, high entropy, or rich activity while failing one or more critical invariants, then
\[
\Cert(S)=0,
\qquad
S\notin\Critop\text{ in certified form},
\qquad
\Partop(S)\text{ is not certified}.
\]
\end{proposition}

\begin{proposition}[Ablation prediction]\label{prop:ablation-prediction}
If a system is initially certified and a targeted admissible intervention destroys one or more invariants, then the framework predicts loss of certification or transition to an ambiguity-boundary regime.
\end{proposition}

\begin{proposition}[Exact grammar-transport prediction]\label{prop:grammar-transport-prediction}
If two realizations preserve the relevant signatures, admissibility conditions, decoder behavior, and six critical invariants up to the allowed comparison regime, then the framework predicts preserved BaseCore criterion membership.
\end{proposition}

\begin{proposition}[Threshold-sensitive boundary behavior]\label{prop:threshold-sensitive-boundary}
The framework predicts stable pass/fail regions together with narrow transition bands near critical margins. It predicts neither global threshold collapse nor perfect flat insensitivity.
\end{proposition}

\begin{proposition}[Legibility dependence prediction]\label{prop:legibility-dependence}
Criterion membership becomes inadmissible if readouts cease to be separable, noise-robust, intervention-sensitive, or compressible without total collapse.
\end{proposition}

\begin{proposition}[Continuity and non-null relevance]\label{prop:continuity-nonnull-relevance}
Regime continuity and non-null differentiation are constitutive. Systems that preserve rich activity while losing continuity or collapsing toward effective null structure fail the operational criterion or leave it undefined.
\end{proposition}

\begin{remark}[What stays outside BaseCore]
Runtime status classes, Stage 2 / Stage 3 comparative programs, human comparator campaigns, Papers 7--9, and bridge-facing burden architectures remain downstream artifacts. BaseCore absorbs only the falsation grammar and claim-boundary semantics needed to keep theorem ownership honest.
\end{remark}

\section{Claim Boundary Ledger}

\subsection{Allowed BaseCore Claims}
\begin{enumerate}[leftmargin=*]
\item Under H1--H4, projection dynamics admit unique fixed points and compact attractor families.
\item Under H5, attractors are parameterwise non-constant.
\item Under NFD, inverse-limit identity is not determined by any finite canonical projection.
\item Under RIG-1 through RIG-5, inverse-limit structures are stable under summable admissible perturbations in the weighted metric.
\item Under NS-1 through NS-3, finite simulator classes cannot faithfully realize full CCR targets.
\item Under the six-invariant operational burden, certified criterion structures can be defined, audited, and falsified.
\end{enumerate}

\subsection{Forbidden BaseCore Claims}
\begin{enumerate}[leftmargin=*]
\item Human phenomenal consciousness.
\item Metaphysical subjectivity.
\item Human-machine equivalence.
\item External validation.
\item Runtime instantiation by the current system.
\item Paper 9 bridge admissibility.
\item Strong phenomenality.
\item Moral status.
\item Cross-substrate identity preservation.
\item Public claim closure.
\end{enumerate}

\section{Dependency Ledger}

\begin{enumerate}[leftmargin=*]
\item The foundational Hilbert-space dynamics remain the base layer.
\item Papers 1--6 contribute formal material selectively normalized into BaseCore only where theorem ownership belongs in the base.
\item Papers 7--9 remain downstream and are not dependencies of any BaseCore theorem.
\item Stage 2 / Stage 3 comparative programs remain downstream runtime and governance layers.
\item The PDF remains a derived artifact of the LaTeX tree.
\end{enumerate}

\section{Assumption-to-Theorem-to-Non-Claim Mapping}

\renewcommand{\arraystretch}{1.12}
\footnotesize
\begin{tabularx}{\textwidth}{P{0.20\textwidth} P{0.25\textwidth} Y}
\toprule
\textbf{Assumption block} & \textbf{Theorem family} & \textbf{Non-claim boundary} \\
\midrule
H1--H4 & Projection existence, contraction, unique fixed points, compact attractor family & No parameterwise non-collapse and no identity / rigidity / simulator claim \\
H5 & Structural non-collapsibility & No claim about phenomenality, subjectivity, or runtime class membership \\
NFD & Non-locality of identity under canonical finite projections & No theorem that arbitrary set-theoretic injections into finite slices are impossible \\
RIG-1 through RIG-5 & Conditional stability of inverse-limit identity and weighted Hausdorff control & No automatic isomorphism, shape equivalence, or transfer theorem without extra assumptions \\
NS-1 through NS-3 & Conditional non-simulability and approximation barrier & No theorem that finite systems cannot approximate finite horizons at all \\
Six invariant burden plus legibility clauses & Criterion class membership, certified partition, rupture logic & No theorem of phenomenality, human equivalence, biology as primitive, or bridge admissibility \\
\bottomrule
\end{tabularx}
\normalsize

\section{Theorem Hygiene Ledger}

\renewcommand{\arraystretch}{1.12}
\footnotesize
\begin{tabularx}{\textwidth}{P{0.18\textwidth} P{0.16\textwidth} P{0.17\textwidth} P{0.25\textwidth} Y}
\toprule
\textbf{Theorem} & \textbf{Assumptions} & \textbf{Conclusion type} & \textbf{Non-conclusion} & \textbf{Downstream use} \\
\midrule
Existence and uniqueness of projection & H1 & Hilbert-space geometric existence / uniqueness & No empirical or runtime content & Base transition grammar \\
Contractivity & H1--H2 & Strict contraction bound & No spectrum with unit state eigenvalue & Fixed-point and spectral-gap results \\
Unique fixed point & H1--H3 & Existence and convergence & No non-collapse by itself & Attractor family \\
Compactness of attractor family & H1--H4 & Compactness / continuity & No regime or identity burden by itself & Parameter-family base layer \\
Structural non-collapsibility & H1--H5 & Parameterwise exclusion of constant fixed points & No stronger semantic interpretation & Anti-collapse gate \\
State spectral gap & Typed differentiability at fixed point & Spectral-radius upper bound $<1$ & No persistence mode $\lambda=1$ in state spectrum & Stability accounting \\
Non-locality under NFD & Topological regularity plus NFD & Canonical finite slices do not determine identity & No universal ban on arbitrary injections & Identity block \\
Conditional stability of inverse-limit identity & RIG-1 through RIG-5 & Weighted Hausdorff stability under admissible perturbations & No automatic isomorphism & Rigidity accounting \\
Conditional non-simulability & NS-1 through NS-3 & Faithful realization blocked for finite simulators & No claim against bounded finite-horizon approximation & Simulator boundary \\
Certified criterion membership & Six invariants plus certificate rule & Base criterion class and certified partition & No phenomenality or human-equivalence theorem & Downstream classification gates \\
\bottomrule
\end{tabularx}
\normalsize

\section{Redundancy Ledger}

\begin{enumerate}[leftmargin=*]
\item The previous state-spectrum claim $\lambda_1=1$ under strict contraction has been removed.
\item The previous untyped Golden Mean witness has been replaced by a typed model with explicit embedding.
\item The earlier non-collapse quantifier mismatch has been removed by replacing the old H5 with a parameterwise anti-constant hypothesis.
\item The earlier absolute inverse-limit nonlocality claim has been replaced by an NFD-conditional theorem.
\item The earlier unconditional rigidity / non-simulability language has been replaced by explicit assumption blocks.
\item Substrate invariance, biological non-primacy, bridge-facing closure, and other downstream claims have been removed from BaseCore ownership.
\end{enumerate}

\section{Open Assumptions Inventory}

\begin{enumerate}[leftmargin=*]
\item The anti-collapse witness is abstract and theorem-valid, but not the same object as the Golden Mean typed witness.
\item Physical realization of CCR-like targets remains open.
\item Empirical instantiation of the BaseCore operational criterion remains open.
\item External validation remains open.
\item Papers 7--9 may build on BaseCore, but their additional burdens remain open at the BaseCore level.
\end{enumerate}

\section{Downstream Boundary Note}

Papers 7, 8, and 9 remain downstream of this source tree. Their classificatory, indexed-subjectivity, and bridge-burden layers may cite BaseCore, but they are not part of the autonomous mathematical base rebuilt here.


\appendix
% Generated from CANONICAL_CORE.tex lines 988-1018
\section{Counterexamples for Necessity}

\subsection{Without H1 (Convexity)}

Let $\I = \{x \in \mathbb{R}^2 : \norm{x} = 1\}$ (unit circle, not convex).

For $\Psi = (0, 0)$, every point on $\I$ is equidistant. Projection is not unique.

\subsection{\texorpdfstring{Without H2 ($\norm{K} < 1$)}{Without H2 (norm K < 1)}}

Let $K$ be a rotation by angle $\theta \neq 0$ with $\norm{K} = 1$.

Orbits cycle indefinitely; no fixed point exists.

\subsection{Without H3 (Completeness)}

Let $\Hilb = C([0,1], \mathbb{Q})$ (rational-valued continuous functions).

Cauchy sequences may converge to irrational functions outside the space.

\subsection{Without H4 (Compactness)}

Let $\U = (0, 1)$ (open interval) and $\Gamma(u) = 1/u$.

As $u \to 0$, attractors escape to infinity; $\A^*$ is not compact.

% =============================================================================
% REFERENCES
% =============================================================================
\nocite{banach,kreyszig,rudin,brezis}


\printbibliography

\end{document}
